\BourbakiAppendixExercises{Exercises}

\begin{enumerate}
\def\labelenumi{\arabic{enumi})}
\tightlist
\item
  Let \(\mathrm{A}\) be a \(k\) -algebra.
\end{enumerate}

\begin{enumerate}
\def\labelenumi{\alph{enumi})}
\item
  Prove that the only regular (left or right) ideal containing AA is A.
\item
  A two-sided ideal \(\mathfrak{a}\) of \(A\) is regular as a left ideal and as a right ideal if and only if the \(k\) -algebra \(A / \mathfrak{a}\) is unital.
\end{enumerate}

\begin{enumerate}
\def\labelenumi{\arabic{enumi})}
\setcounter{enumi}{1}
\tightlist
\item
  Let \(\mathbf{A}\) be a \(k\) -algebra, and let \(\mathfrak{a}\) and \(\mathfrak{b}\) be regular (left) ideals of \(\mathbf{A}\) .
\end{enumerate}

\begin{enumerate}
\def\labelenumi{\alph{enumi})}
\item
  The ideal \(a \cap b\) is regular if and only if there exist right units u modulo a and v modulo b such that \(u - v \in a + b\) .
\item
  Deduce that \(a \cap b\) is regular if a is two-sided and right regular, or if a is maximal. c) Let K be a commutative field, L be the free associative K-algebra in two variables X and Y (III, §2, No.~7, p.~448), and A be the K-subalgebra of L consisting of the elements of degree \(\geqslant 1\) . Prove that the ideals \(A(1 - X)\) and \(A(1 - Y)\) of A are regular but that their intersection is reduced to 0.
\end{enumerate}

\begin{enumerate}
\def\labelenumi{\arabic{enumi})}
\setcounter{enumi}{2}
\tightlist
\item
  Let \(\mathbf{A}\) be a \(k\) -algebra, \(\mathfrak{a}\) a regular left ideal of \(\mathbf{A}\) , and \(u\) and \(v\) right units modulo \(\mathfrak{a}\) . \(a)\) Give an example where \(u - v\) does not belong to \(\mathfrak{a}\) (cf.~III, §9, p.~644, Exercise 1). Deduce an example where the mapping \(\widetilde{\mathfrak{a}} \mapsto \widetilde{\mathfrak{a}} \cap \mathbf{A}\) (Proposition 3 of VIII, p.~437) is not injective.
\end{enumerate}

\begin{enumerate}
\def\labelenumi{\alph{enumi})}
\setcounter{enumi}{1}
\tightlist
\item
  Suppose that the ideal \(\mathfrak{a}\) is two-sided and that the \(k\) -algebra \(B = A / \mathfrak{a}\) does not contain any nonzero two-sided ideal \(\mathfrak{b}\) such that \(B\mathfrak{b} = 0\) . Prove that \(u - v\) belongs to \(\mathfrak{a}\) . This is, in particular, the case when \(A\) is commutative.
\end{enumerate}

\begin{enumerate}
\def\labelenumi{\arabic{enumi})}
\setcounter{enumi}{3}
\tightlist
\item
  Let A be a \(k\) -algebra and \(\mathfrak{a}\) a left ideal of A; we view \(\mathfrak{a}\) as a \(k\) -algebra. Prove that we have \(\Re(A) \cap \mathfrak{a} \subset \Re(\mathfrak{a})\) (apply Proposition 7 of VIII, p.~441) and that we have equality when \(\mathfrak{a}\) is two-sided (observe that every simple pseudomodule over A is simple over \(\mathfrak{a}\) or annihilated by \(\mathfrak{a}\) ). Give an example of a left ideal \(\mathfrak{a}\) such that \(\Re(A) \cap \mathfrak{a} \neq \Re(\mathfrak{a})\) (take A to be a matrix algebra).
\end{enumerate}

¶ 5) Suppose that every pseudomodule M over A is the direct sum of AM and the sub-pseudomodule of M consisting of the elements annihilated by A. Prove that A is unital. (Taking \(M = \widetilde{A}_{s}\) , prove that A admits a right unit element u. Taking \(M = A_{s}\) , prove that we have \(Ax \neq 0\) for every nonzero element x of A. Conclude that u is also a left unit element.)

\begin{enumerate}
\def\labelenumi{\arabic{enumi})}
\setcounter{enumi}{5}
\tightlist
\item
  We say that a k-algebra is primitive if it admits a pseudomodule that is simple and faithful (that is, whose annihilator is reduced to 0).
\end{enumerate}

\begin{enumerate}
\def\labelenumi{\alph{enumi})}
\item
  Show that if the algebra \(\tilde{\mathbf{A}}\) is primitive, then so is \(\mathbf{A}\) .
\item
  If the algebra A is primitive and unital, then the algebra \(\widetilde{A}\) is not primitive.
\item
  Suppose that k is a field and that the algebra A is primitive and not unital. Prove that \(\widetilde{A}\) is primitive (prove that if there exists a faithful pseudomodule over A that is not faithful as an \(\widetilde{A}\) -module, then A admits a right unit element u, and deduce a contradiction).
\end{enumerate}

\begin{enumerate}
\def\labelenumi{\arabic{enumi})}
\setcounter{enumi}{6}
\item
  We say that the \(k\) -algebra A is Artinian if every decreasing sequence of left ideals of A is stationary. The ring \(\widetilde{\mathrm{A}}\) is Artinian if and only if the algebra A and the ring \(k\) are Artinian. Prove that if \(k\) is a field, then an Artinian algebra A without radical is unital (observe that the algebra \(\widetilde{\mathrm{A}}\) is semisimple).
\item
  Suppose that \(k\) is a field and that \(A\) is a primitive algebra over \(k\) . Prove that if \(A\) admits a finite-dimensional nonzero left ideal over \(k\) , then the algebra \(A\) is finite-dimensional over \(k\) and consequently a simple ring (cf.~Exercise 7).
\item
  Suppose that k is a field and that A is a nilpotent k-algebra, that is, that there exists an integer n such that \(A^{n}=0\) . Prove that A is finite-dimensional over k (use induction on the least integer n such that \(A^{n}=0\) ).
\item
  Suppose that \(k\) is a field, that \(A\) is an infinite-dimensional \(k\) -algebra, and that every proper left ideal of \(A\) is finite-dimensional over \(k\) .
\end{enumerate}

\begin{enumerate}
\def\labelenumi{\alph{enumi})}
\item
  Prove that we have \(Aa = 0\) for every left ideal \(a \neq A\) (consider the homomorphism \(\gamma : A \to \operatorname{End}_{k}(a)\) ).
\item
  Let \(\mathfrak{r}\) be the sum of the maximal left ideals of A. Prove that \(\mathfrak{r}\) is different from A (use Exercise 9), that \(\mathfrak{r}\) is a two-sided ideal, and that \(A / \mathfrak{r}\) is a field (cf.~Exercise 7). c) Let \(u\) be a (left or right) unit modulo \(\mathfrak{r}\) . Prove that we have \(\mathfrak{ru} = \mathfrak{r}\) (observe that we have \(Au = A\) ).
\item
  Prove that rA is zero (consider the right annihilator of r in A and reason as in a)) and then that r is zero. Conclude that A is a field.
\end{enumerate}

¶ 11) Suppose that k is a field, that A is a k-algebra, and that every decreasing sequence of subalgebras of A is stationary.

\begin{enumerate}
\def\labelenumi{\alph{enumi})}
\item
  Prove that the algebra \(\mathbf{A}\) is Artinian and that its radical is finite-dimensional over \(k\) (apply Exercise 9).
\item
  Let S be a simple component of A/ \(\Re(A)\) (Exercise 7); it is isomorphic to a matrix algebra over a field D, with center Z. Prove that D is finite-dimensional over Z (if D contains an element that is transcendental over Z, construct an infinite strictly decreasing sequence of subfields of D; if D is algebraic and infinite-dimensional over Z, construct an infinite strictly increasing sequence of subfields of D, and consider their commutants in D).
\item
  Prove that if Z is infinite-dimensional over k, then we have S = D (consider the nilpotent subalgebras of B).
\item
  Prove that Z is an algebraic extension of k. Let p be the characteristic exponent of k and L the largest separable extension of k contained in Z. Prove that \(L \cap Z^{p}\) has finite degree over \(L^{p}\) (consider a p-basis of \(L \cap Z^{p}\) over \(L^{p}\) ).
\item
  Every decreasing sequence of subextensions of L is stationary. Give an example to show that this condition does not imply that L has finite degree over k (cf.~V, §12, p.~167, Exercise 3).
\item
  Conversely, let B be a k-algebra satisfying the conditions of a) such that every simple component of \(\mathrm{B}/\Re(\mathrm{B})\) satisfies the conditions of b), c), d), and e). Prove that every decreasing sequence of subalgebras of B is stationary.
\end{enumerate}

(Reduce to the case when B is a field, with center Z. Let \((\mathrm{D}_{n})\) be a decreasing sequence of subfields of B; we may assume that the field D generated by \(D_{n}\) and Z is independent of n.~Observe that we can write \(D = E \otimes_{L} Z\) , where L is a subfield of Z of finite degree over k and E is a field of finite degree over L whose center contains L. Prove that \(D_{n}\) is then the field generated by E and \(D_{n} \cap Z\) . We are thus reduced to the case when B is a commutative field that is a p-radical extension of k. If \((\mathrm{F}_{n})\) is a decreasing sequence of subextensions of B, consider the subfield \(k(\mathrm{F}_{n}^{p^{\infty}})\) , where \(F_{n}^{p^{\infty}} = \cap_{f} F_{n}^{p^{f}}\) .)

¶ 12) a) Suppose that k is a field, that the k-algebra A is Artinian, and that every increasing sequence of subalgebras of A is stationary. Prove that A has finite degree over k. (Reason as in Exercise 11. Note that the field \(k(\mathrm{X})\) contains an infinite strictly increasing sequence of subalgebras, using VII, §1, No.~5, p.~6, Proposition 6.) b) Prove that in the polynomial algebra \(k[X]\) , every increasing sequence of subalgebras is stationary. (Let \(f \in k[X]\) be a nonconstant polynomial and A be the k-subalgebra of \(k[X]\) generated by 1 and f.~Observe that A is a principal ideal domain and \(k[X]\) is a finitely generated A-module.)

\begin{enumerate}
\def\labelenumi{\arabic{enumi})}
\setcounter{enumi}{12}
\item
  Suppose that \(k\) is a field and that \(A\) admits a finite basis over \(k\) consisting of nilpotent elements. Prove that \(A\) is nilpotent (Exercise 9; reduce to the case when the algebra \(A\) is semisimple (cf.~Exercise 7), and then of the form \(\mathbf{M}_n(\mathrm{K})\) by extending the base field).
\item
  Let A and B be k-algebras; an (A,B)-pseudobimodule is a set E endowed with the structure of a left pseudomodule over A and the structure of a right pseudomodule over B such that the underlying k-module structures coincide. This corresponds to endowing E with the structure of a bimodule over the algebras \(\widetilde{A}\) and \(\widetilde{B}\) . A multiplier of E is a pair (S,D), where S:B→E is a homomorphism of right B-pseudomodules and D:A→E is a homomorphism of left A-pseudomodules, satisfying \(a\mathrm{S}(b)=\mathrm{D}(a)b\) for \(a\in A\) , \(b\in B\) . We denote by \(\mathcal{M}(\mathrm{E})\) the set of multipliers of E; it admits a natural k-module structure.
\end{enumerate}

\(a\) ) For every \(x\in \mathrm{E}\) , we denote by \(\mathrm{D}_x\) (resp. \(\mathrm{S}_x\) ) the mapping \(a\mapsto ax\) (resp. \(b\mapsto xb\) ). Show that the pair \(m_{x} = (\mathrm{D}_{x},\mathrm{S}_{x})\) is a multiplier of E.

\begin{enumerate}
\def\labelenumi{\alph{enumi})}
\setcounter{enumi}{1}
\tightlist
\item
  For every \(m = (\mathrm{S}, \mathrm{D}) \in \mathcal{M}(\mathrm{E})\) , \(a \in A\) , and \(b \in B\) , set \(am = m_{\mathrm{D}(a)}\) and \(mb = m_{\mathrm{S}(b)}\) . Prove that this defines on the k-module \(\mathcal{M}(\mathrm{E})\) the structure of a (A, B)-pseudobimodule and that the mapping \(x \mapsto m_x\) is a homomorphism of (A, B)-pseudobimodules.
\end{enumerate}

\begin{enumerate}
\def\labelenumi{\arabic{enumi})}
\setcounter{enumi}{14}
\tightlist
\item
  Let \(\mathrm{A}\) be a \(k\) -algebra. A multiplier of \(\mathrm{A}\) is a multiplier of the (A,A)-pseudobimodule \(\mathrm{A}\) . We denote the \(k\) -algebra of multipliers of \(\mathrm{A}\) by \(\mathcal{M}(\mathrm{A})\) .
\end{enumerate}

\begin{enumerate}
\def\labelenumi{\alph{enumi})}
\item
  Prove that the mapping \(\mu: a \mapsto m_{a}\) is a k-algebra homomorphism and that we have \(am = m_{a}m\) and \(ma = mm_{a}\) for \(a \in A, m \in \mathcal{M}(A)\) . The image of the homomorphism \(\mu\) is therefore a two-sided ideal of the ring \(\mathcal{M}(A)\) .
\item
  A two-sided ideal b of a k-algebra A is called faithful if its left and right annihilators are zero. Show that if A is faithful in A, then the homomorphism \(\mu\) is injective and the ideal \(\mu(\mathrm{A})\) is faithful in \(\mathcal{M}(\mathrm{A})\) .
\item
  Let B be a unital k-algebra containing A as a two-sided ideal. For \(b \in B\) , let \(\sigma_{b}\) and \(\delta_{b}\) be the mappings from A to A such that for \(a \in A\) , we have \(\sigma_{b}(a) = ba\) and \(\delta_{b}(a) = ab\) , and set \(\pi(b) = (\sigma_{b}, \delta_{b})\) . Show that \(\pi(b) \in \mathcal{M}(A)\) for every \(b \in B\) and that the resulting mapping \(\pi : B \to \mathcal{M}(A)\) is a ring homomorphism whose restriction to A is the mapping \(a \mapsto m_{a}\) . Show that \(\pi\) is injective if the ideal A is faithful in B.
\end{enumerate}

\chapter{Determinants over a Noncommutative Field}

In this appendix, D is a field, and \(D_{ab}^{*}\) is the quotient of the multiplicative group \(D^{*}\) of D by its derived group (I, §6, No.~2, p.~70). The group \(D_{ab}^{*}\) is commutative, and \(D_{ab}^{*}\) can be identified with \(D^{*}\) if the field D is commutative. We denote the canonical homomorphism from \(D^{*}\) to \(D_{ab}^{*}\) by \(\pi\) .

\section{A Generalization of Alternating Multilinear Forms}

Let \(\mathrm{V}\) be a right vector space over the field \(\mathrm{D}\) , of finite dimension \(n \geqslant 0\) . We denote by \(\mathrm{B(V)}\) the set of bases of \(\mathrm{V}\) and by \(\Omega(\mathrm{V})\) the set of mappings \(\omega\) from \(\mathrm{B(V)}\) to \(\mathrm{D}_{\mathrm{ab}}^{*}\) that satisfy the following two conditions:

\begin{enumerate}
\def\labelenumi{\alph{enumi})}
\tightlist
\item
  If \(\lambda_1, \ldots, \lambda_n\) are the elements of \(D^*\) , then we have
\end{enumerate}

\[
\omega (v _ {1} \lambda_ {1}, \ldots , v _ {n} \lambda_ {n}) = \pi (\lambda_ {1} \dots \lambda_ {n}) \omega (v _ {1}, \ldots , v _ {n})\tag{1}
\]

for every basis \((v_{1},\ldots ,v_{n})\) of V.

\begin{enumerate}
\def\labelenumi{\alph{enumi})}
\setcounter{enumi}{1}
\tightlist
\item
  If \(i\) and \(j\) are two distinct integers in the interval \([1, n]\) , then we have
\end{enumerate}

\[
\omega (v _ {1}, \ldots , v _ {i - 1}, v _ {i} + v _ {j}, v _ {i + 1}, \ldots , v _ {n}) = \omega (v _ {1}, \ldots , v _ {n})\tag{2}
\]

for every basis \((v_{1},\ldots ,v_{n})\) of V.

Let \(\omega\) be an element of \(\Omega(\mathrm{V})\) . Let \((v_1, \ldots, v_n)\) be a basis of \(\mathrm{V}\) and \(i\) an integer in the interval \([1, n-1]\) ; by property b), we have

\[
\begin{array}{r l} & {\omega (v _ {1}, \ldots , v _ {i}, v _ {i + 1}, \ldots , v _ {n}) = \omega (v _ {1}, \ldots , v _ {i} + v _ {i + 1}, v _ {i + 1}, \ldots , v _ {n})} \\ & {\qquad = \omega (v _ {1}, \ldots , v _ {i} + v _ {i + 1}, v _ {i + 1} - (v _ {i} + v _ {i + 1}), \ldots , v _ {n})} \\ & {\qquad = \omega (v _ {1}, \ldots , v _ {i} + v _ {i + 1}, - v _ {i}, \ldots , v _ {n})} \\ & {\qquad = \omega (v _ {1}, \ldots , (v _ {i} + v _ {i + 1}) - v _ {i}, - v _ {i}, \ldots , v _ {n}),} \end{array}
\]

and therefore

\[
\omega (v _ {1}, \dots , v _ {i}, v _ {i + 1}, \dots , v _ {n}) = \pi (- 1) \omega (v _ {1}, \dots , v _ {i + 1}, v _ {i}, \dots , v _ {n}).\tag{3}
\]

Since the symmetric group \(\mathfrak{S}_n\) is generated by the transpositions of two consecutive elements of the interval \([1,n]\) (I, §5, No.~7, p.~63, Proposition 9), formula (3) generalizes to

\[
\omega (v _ {\sigma (1)}, \ldots , v _ {\sigma (n)}) = \pi (\varepsilon (\sigma)) \omega (v _ {1}, \ldots , v _ {n}),\tag{4}
\]

where \(\sigma\) belongs to \(\mathfrak{S}_n\) and \(\varepsilon (\sigma)\) is its signature (I, §5, No.~7, p.~64).

Formula (2) generalizes as follows. First, for every \(\lambda\) in \(\mathrm{D}^*\) , we have

\[
\begin{array}{r l} & {\omega (v _ {1}, \ldots , v _ {n}) = \pi (\lambda) \omega (v _ {1}, \ldots , v _ {i} \lambda^ {- 1}, \ldots , v _ {n})} \\ & {\qquad = \pi (\lambda) \omega (v _ {1}, \ldots , v _ {i} \lambda^ {- 1} + v _ {j}, \ldots , v _ {n})} \\ & {\qquad = \omega (v _ {1}, \ldots , (v _ {i} \lambda^ {- 1} + v _ {j}) \lambda , \ldots , v _ {n}),} \end{array}
\]

that is,

\[
\omega (v _ {1}, \dots , v _ {n}) = \omega (v _ {1}, \dots , v _ {i} + v _ {j} \lambda , \dots , v _ {n}).\tag{5}
\]

By an immediate induction, we deduce

\[
\omega (v _ {1}, \ldots , v _ {n}) = \omega (v _ {1}, \ldots , v _ {i} + w, \ldots , v _ {n})\tag{6}
\]

for every linear combination \(w = \sum_{j \neq i} v_j \lambda_j\) of the vectors \(v_j\) for j different from i in the interval \([1, n]\) .

\section{A Uniqueness Theorem}

In this subsection, for any strictly positive integer m, any integer i in the interval \([1,m]\) , and any sequence \((v_{1},\ldots,v_{m})\) in \(V^{m}\) , we write \((v_{1},\ldots,\widehat{v_{i}},\ldots,v_{m})\) for the sequence \((v_{1},\ldots,v_{i-1},v_{i+1},\ldots,v_{m})\) in \(V^{m-1}\) .

PROPOSITION 1. --- Let W be a hyperplane in V and e be a vector in V---W. Let \(\varphi\) be an element of \(\Omega(\mathrm{W})\) . There exists a unique element \(\omega\) of \(\Omega(\mathrm{V})\) such that we have

\[
\omega (w _ {1}, \ldots , w _ {n - 1}, e) = \varphi (w _ {1}, \ldots , w _ {n - 1})\tag{7}
\]

for every basis \((w_{1},\ldots ,w_{n - 1})\) of W.

\begin{enumerate}
\def\labelenumi{\Alph{enumi})}
\tightlist
\item
  Uniqueness of \(\omega\) : Let \(\omega\) be an element of \(\Omega(\mathrm{V})\) . Let \((v_{1},\ldots,v_{n})\) be a basis of V, and let \(\mu_{1},\ldots,\mu_{n}\) be the coordinates of e in this basis. Denote by I the set of integers from 1 to n (inclusive) such that \(\mu_{i}\neq0\) ; it is not empty. Let i be an element of I. The sequence \((v_{1},\ldots,\widehat{v_{i}},\ldots,v_{n},e)\) is a basis of V, and since \(e-v_{i}\mu_{i}\) is a linear combination of the sequence \((v_{1},\ldots,\widehat{v_{i}},\ldots,v_{n})\) , formula (6) implies
\end{enumerate}

\[
\omega (v _ {1}, \ldots , \widehat {v _ {i}}, \ldots , v _ {n}, v _ {i} \mu_ {i}) = \omega (v _ {1}, \ldots , \widehat {v _ {i}}, \ldots , v _ {n}, e).\tag{8}
\]

It follows from formulas (1) and (3) that the left-hand side of this equality is equal to \(\pi(-1)^{n-i}\pi(\mu_i)\omega(v_1,\ldots,v_n)\) . Denote by p the projector of V with image W and kernel eD. The vectors \(v_j - p(v_j)\) are proportional to e, and repeatedly applying formula (6) shows that the right-hand side of formula (8) is equal to \(\omega(p(v_1),\ldots,\widehat{p(v_i)},\ldots,p(v_n),e)\) .

It follows from this that if \(\omega\) satisfies relation (7), then we have

\[
\omega (v _ {1}, \dots , v _ {n}) = \pi (- 1) ^ {n - i} \pi (\mu_ {i}) ^ {- 1} \varphi (p (v _ {1}), \dots , \widehat {p (v _ {i})}, \dots , p (v _ {n}))\tag{9}
\]

for every i in I, which proves the uniqueness of \(\omega\) .

\begin{enumerate}
\def\labelenumi{\Alph{enumi})}
\setcounter{enumi}{1}
\tightlist
\item
  Construction of \(\omega\) : Let \((v_{1},\ldots,v_{n})\) be a basis of V; we define \(\mu_{1},\ldots,\mu_{n}\) and I as above. For every \(i\) in I, the assumption \(\mu_{i}\neq0\) implies that the sequence \((v_{1},\ldots,\widehat{v_{i}},\ldots,v_{n},e)\) generates the space V, so the sequence \((p(v_{1}),\ldots,\widehat{p(v_{i})},\ldots,p(v_{n}))\) generates the space W = p(V). In other words, this last sequence is a basis of W, and
\end{enumerate}

\[
t _ {i} = \pi (- 1) ^ {n - i} \pi (\mu_ {i}) ^ {- 1} \varphi (p (v _ {1}), \ldots , \widehat {p (v _ {i})}, \ldots , p (v _ {n}))\tag{10}
\]

defines an element of \(\mathrm{D}_{\mathrm{ab}}^{*}\) .

Let i and j be two elements of I such that i \textless{} j; let us prove the equality \(t_{i} = t_{j}\) . By definition, the vector \(v_{i}\mu_{i} + v_{j}\mu_{j} - e\) is a linear combination of the vectors \(p(v_{k})\) for k different from i and j. Consequently, \(p(v_{i})\mu_{i} + p(v_{j})\mu_{j}\) is a linear combination of the vectors \(p(v_{k})\) for k different from i and j. By formula (6), we therefore have

\[
\begin{array}{c} \varphi (p (v _ {1}), \ldots , \widehat {p (v _ {i})}, \ldots , \widehat {p (v _ {j})}, \ldots , p (v _ {n}), p (v _ {i}) \mu_ {i}) \\ = \varphi (p (v _ {1}), \ldots , \widehat {p (v _ {i})}, \ldots , \widehat {p (v _ {j})}, \ldots , p (v _ {n}), - p (v _ {j}) \mu_ {j}). \end{array}
\]

By applying formulas (1) and (3), we deduce

\[
\begin{array}{c} \varphi (p (v _ {1}), \ldots , \widehat {p (v _ {j})}, \ldots , p (v _ {n})) \pi (\mu_ {i}) \pi (- 1) ^ {n - i - 1} \\ = \varphi (p (v _ {1}), \ldots , \widehat {p (v _ {i})}, \ldots , p (v _ {n})) \pi (- \mu_ {j}) \pi (- 1) ^ {n - j}, \end{array}
\]

and therefore \(t_i = t_j\) .

Having proved this, we define \(\omega(v_1, \ldots, v_n)\) as the common value of the \(t_i\) for \(i\) in I. We have thus constructed a mapping \(\omega\) from B(V) to D \(_{\text{ab}}^*\) that satisfies relation (9). Let \((w_1, \ldots, w_{n-1})\) be a basis of W; if we set \(v_i = w_i\) for \(1 \leqslant i \leqslant n-1\) and \(v_n = e\) , then we have I = \{n\} and \(\mu_n = 1\) . We moreover have \(p(v_i) = w_i\) for \(1 \leqslant i \leqslant n-1\) , and the formula \(\omega(w_1, \ldots, w_{n-1}, e) = \varphi(w_1, \ldots, w_{n-1})\) is a specific case of (9).

We must now prove that \(\omega\) belongs to \(\Omega (\mathrm{V})\) .

\begin{enumerate}
\def\labelenumi{\Alph{enumi})}
\setcounter{enumi}{2}
\tightlist
\item
  Proof of formula (1): Let \(\lambda_1, \ldots, \lambda_n\) be elements of \(D^*\) and \((v_1, \ldots, v_n)\) be a basis of \(V\) . We define \(\mu_1, \ldots, \mu_n\) and \(I\) as above. Choose an \(i\) in \(I\) . Since we have \(e = \sum_{j=1}^{n} (v_j \lambda_j)(\lambda_j^{-1} \mu_j)\) , formula (9) implies
\end{enumerate}

\[
\begin{array}{l l} (1 1) & \omega (v _ {1} \lambda_ {1}, \ldots , v _ {n} \lambda_ {n}) \\ & \qquad = \pi (- 1) ^ {n - i} \pi (\lambda_ {i} ^ {- 1} \mu_ {i}) ^ {- 1} \varphi (p (v _ {1}) \lambda_ {1}, \ldots , \widehat {p (v _ {i}) \lambda_ {i}}, \ldots , p (v _ {n}) \lambda_ {n}). \end{array}
\]

But we have \(\pi (\lambda_i^{-1}\mu_i)^{-1} = \pi (\mu_i)^{-1}\pi (\lambda_i)\) and

\[
\begin{array}{c} \varphi (p (v _ {1}) \lambda_ {1}, \ldots , \widehat {p (v _ {i}) \lambda_ {i}}, \ldots , p (v _ {n}) \lambda_ {n}) \\ = \varphi (p (v _ {1}), \ldots , \widehat {p (v _ {i})}, \ldots , p (v _ {n})) \pi (\lambda_ {1} \ldots \widehat {\lambda_ {i}} \ldots \lambda_ {n}). \end{array}
\]

Comparing formulas (9) and (11) leads to the desired relation

\[
\omega (v _ {1} \lambda_ {1}, \ldots , v _ {n} \lambda_ {n}) = \omega (v _ {1}, \ldots , v _ {n}) \pi (\lambda_ {1} \ldots \lambda_ {n}).
\]

\begin{enumerate}
\def\labelenumi{\Alph{enumi})}
\setcounter{enumi}{3}
\tightlist
\item
  Proof of the formula (2): Let \((v_{1},\ldots,v_{n})\) be a basis of V and i,j be two distinct integer in the interval {[}1,n{]}; we define \(\mu_{1},\ldots,\mu_{n}\) as before. Consider the basis \((v_{1}^{\prime},\ldots,v_{n}^{\prime})\) of V defined by \(v_{i}^{\prime}=v_{i}+v_{j}\) and \(v_{k}^{\prime}=v_{k}\) for \(k\neq i\) , and let us introduce the coordinates \(\mu_{1}^{\prime},\ldots,\mu_{n}^{\prime}\) of e with respect to this basis; they satisfy \(\mu_{j}^{\prime}=\mu_{j}-\mu_{i}\) and \(\mu_{k}^{\prime}=\mu_{k}\) for \(k\neq j\) . We set \(t=\omega(v_{1},\ldots,v_{n})\) and \(t^{\prime}=\omega(v_{1}^{\prime},\ldots,v_{n}^{\prime})\) . We must prove that t and \(t^{\prime}\) are equal.
\end{enumerate}

Let us first note that, by the definition of \(\omega\) , we have

\begin{enumerate}
\def\labelenumi{(\arabic{enumi})}
\setcounter{enumi}{11}
\tightlist
\item
\end{enumerate}

\[
t = \pi (- 1) ^ {n - k} \pi (\mu_ {k}) ^ {- 1} \varphi (p (v _ {1}), \ldots , \widehat {p (v _ {k})}, \ldots , p (v _ {n})),\tag{13}
\]

\[
t ^ {\prime} = \pi (- 1) ^ {n - k} \pi (\mu_ {k} ^ {\prime}) ^ {- 1} \varphi (p (v _ {1} ^ {\prime}), \ldots , \widehat {p (v _ {k} ^ {\prime})}, \ldots , p (v _ {n} ^ {\prime}))
\]

for every \(k\) such that \(\mu_{k}\) and \(\mu_k^\prime\) are nonzero.

\begin{enumerate}
\def\labelenumi{\alph{enumi})}
\item
  If \(\mu_i \neq 0\) , then the sequences \((p(v_1), \ldots, \widehat{p(v_k)}, \ldots, p(v_n))\) and \((p(v_1'), \ldots, \widehat{p(v_k'}, \ldots, p(v_n'))\) are equal, and we have \(\mu_i = \mu_i'\) , and therefore \(t = t'\) .
\item
  If there exists an index \(k\) different from \(i\) and \(j\) such that \(\mu_k \neq 0\) , then we have \(p(v_l') = p(v_l)\) for \(l \neq i\) and \(p(v_i') = p(v_i) + p(v_j)\) . The elements \(p(v_i)\)
\end{enumerate}

and \(p(v_{j})\) both belong to the sequence \((p(v_{1}),\ldots ,\widehat{p(v_{k})},\ldots ,p(v_{n}))\) . Since \(\varphi\) belongs to \(\Omega (\mathrm{W})\) , formula (2) of VIII, p.~447 applies; it gives

\[
\varphi (p (v _ {1}), \dots , \widehat {p (v _ {k})}, \dots , p (v _ {n})) = \varphi (p (v _ {1} ^ {\prime}), \dots , \widehat {p (v _ {k} ^ {\prime})}, \dots , p (v _ {n} ^ {\prime})).
\]

But we also have \(\mu_{k} = \mu_{k}^{\prime}\) , and therefore \(t = t'\) .

\begin{enumerate}
\def\labelenumi{\alph{enumi})}
\setcounter{enumi}{2}
\tightlist
\item
  It remains to examine the case when the only index \(k\) such that \(\mu_k \neq 0\) is \(j\) . We then have \(\mu_j' = \mu_j\) , \(e = v_j\mu_j\) , and \(p(v_j) = 0\) . Set \(k = j\) in formulas (12) and (13). Since the sequences \((p(v_1), \ldots, \widehat{p(v_j)}, \ldots, p(v_n))\) and \((p(v_1'), \ldots, \widehat{p(v_j'), \ldots, p(v_n'))}\) are equal, we have \(t = t'\) .
\end{enumerate}

COROLLARY 1. --- Let \((e_1, \ldots, e_n)\) be a basis of V and t an element of \(\mathrm{D}_{\mathrm{ab}}^*\) . There exists a unique element \(\omega\) of \(\Omega(\mathrm{V})\) such that \(\omega(e_1, \ldots, e_n) = t\) .

Let us prove the corollary by induction on n.~If n = 0, then the only basis of V is the empty one, so the assertion is true. Now suppose \(n \geqslant 1\) ; denote by W the subspace of V generated by \(e_{1}, \ldots, e_{n-1}\) . By Proposition 1, there exists a bijection \(\Lambda\) from \(\Omega(\mathrm{V})\) to \(\Omega(\mathrm{W})\) satisfying

\[
(\Lambda (\omega)) (w _ {1}, \ldots , w _ {n - 1}) = \omega (w _ {1}, \ldots , w _ {n - 1}, e _ {n})
\]

for every basis \((w_{1},\ldots,w_{n-1})\) of W and every \(\omega\) in \(\Omega(\mathrm{V})\) . By the induction hypothesis, there exists a unique element \(\varphi\) of \(\Omega(\mathrm{W})\) such that \(\varphi(e_{1},\ldots,e_{n-1})=t\) . The relation \(\omega(e_{1},\ldots,e_{n})=t\) , for \(\omega\) in \(\Omega(\mathrm{V})\) , is equivalent to \(\Lambda(\omega)=\varphi\) . Corollary 1 follows.

COROLLARY 2. --- Let \(\omega\) and \(\omega'\) be two elements of \(\Omega(V)\) . There exists a unique element \(t\) of \(\mathrm{D}_{\mathrm{ab}}^{*}\) such that \(\omega' = t\omega\) .

Remark. --- Suppose that the field D is commutative. By definition, B(V) is a subset of \(V^{n}\) . It is clear that the restriction to B(V) of a nonzero alternating n-linear form \(f: V^{n} \to D\) belongs to \(\Omega(V)\) . Moreover, if \((e_{1}, \ldots, e_{n})\) is a basis of V and t is a nonzero element of D, then there exists a unique alternating n-linear form f such that \(f(e_{1}, \ldots, e_{n}) = t\) (III, §7, No.~4, p.~511 and III, §7, No.~8, p.~518). By Corollary 1, the set \(\Omega(V)\) consists of the restrictions to B(V) of the nonzero alternating n-linear forms.

\section{Determinant of an Automorphism}

We first suppose that the field D is commutative. In view of the remark above, the determinant of an automorphism u of V is the unique element \(\det u\) of \(D^{*}\) such that we have

\[
\omega (u (v _ {1}), \ldots , u (v _ {n})) = (\det u) \omega (v _ {1}, \ldots , v _ {n})\tag{14}
\]

for every basis \((v_{1},\ldots ,v_{n})\) of V and every element \(\omega\) of \(\Omega (\mathrm{V})\) .

Let us return to the case when D is no longer supposed to be commutative.

PROPOSITION 2. --- a) Let u be an automorphism of V. There exists a unique element of \(D_{ab}^{*}\) , denoted by det u and called the determinant of u, such that we have

\[
\omega (u (v _ {1}), \ldots , u (v _ {n})) = (\det u) \omega (v _ {1}, \ldots , v _ {n})\tag{15}
\]

for every basis \((v_{1},\ldots ,v_{n})\) of V and every \(\omega\) in \(\Omega (\mathrm{V})\)

\begin{enumerate}
\def\labelenumi{\alph{enumi})}
\setcounter{enumi}{1}
\tightlist
\item
  The mapping \(u \mapsto \det u\) from \(\mathbf{GL}(\mathrm{V})\) to \(\mathrm{D}_{\mathrm{ab}}^{*}\) is a group homomorphism.
\end{enumerate}

Let \(\omega_0\) be an element of \(\Omega (\mathrm{V})\) . The mapping \((v_{1},\ldots ,v_{n})\mapsto\) \(\omega_0(u(v_1),\dots ,u(v_n))\) from B(V) to \(\mathrm{D}_{\mathrm{ab}}^{*}\) belongs to \(\Omega (\mathrm{V})\) . By Corollary 2 of VIII, p.~451, there exists a unique element \(t\) of \(\mathrm{D}_{\mathrm{ab}}^{*}\) such that we have

\[
\omega_ {0} (u (v _ {1}), \ldots , u (v _ {n})) = t \omega_ {0} (v _ {1}, \ldots , v _ {n})
\]

for \((v_{1},\ldots ,v_{n})\) in \(\mathrm{B(V)}\) . If \(\omega\) is another element of \(\Omega (\mathrm{V})\) , then there exists an element \(s\) of \(\mathrm{D}_{\mathrm{ab}}^{*}\) such that

\[
\omega (v _ {1}, \ldots , v _ {n}) = s \omega_ {0} (v _ {1}, \ldots , v _ {n})
\]

for every basis \((v_{1},\ldots ,v_{n})\) of V (loc. cit.). We immediately deduce the relation

\[
\omega (u (v _ {1}), \dots , u (v _ {n})) = t \omega (v _ {1}, \dots , v _ {n}).
\]

This proves a); assertion b) is an immediate consequence of a).

\section{Determinant of a Square Matrix}

Let n be a positive integer. Let us apply the above to the right vector space \(D_{d}^{n}\) over the field D; the elements of \(D_{d}^{n}\) are viewed as matrices with n rows and one column. Let \((\varepsilon_{1},\ldots,\varepsilon_{n})\) be the canonical basis of \(D_{d}^{n}\) . By Corollary 2 of VIII, p.~451, there exists a unique element \(\omega_{0}\) of \(\Omega(\mathrm{D}_{d}^{n})\) such that \(\omega_{0}(\varepsilon_{1},\ldots,\varepsilon_{n})=1\) . If A is an element of \(\mathbf{GL}_{n}(\mathrm{D})\) , then its columns \(a_{1},\ldots,a_{n}\) form a basis of \(D_{d}^{n}\) ; the element \(\omega_{0}(a_{1},\ldots,a_{n})\) of \(D_{ab}^{*}\) is called the determinant of A and denoted by \(\det(A)\) . Since we have \(a_{i}=A\varepsilon_{i}\) for \(1\leqslant i\leqslant n\) , the determinant of A is simply the determinant of the automorphism \(x\mapsto Ax\) of \(D_{d}^{n}\) . In particular, if the field D is commutative, then the determinant of A coincides with that defined in III, §8, No.~3, p.~524.

Let \(\mathrm{V}\) be a right vector space of finite dimension \(n\) over the field \(\mathrm{D}\) , and let \((e_1, \ldots, e_n)\) be a basis of \(\mathrm{V}\) . If \(u\) is an automorphism of \(\mathrm{V}\) and \(A\) is the matrix of \(u\) with respect to the basis \((e_1, \ldots, e_n)\) , then we have \(\det(u) = \det(A)\) .

We denote by \((E_{ij})\) the family of matrix units of \(\mathbf{M}_{n}(\mathrm{D})\) (II, §10, No.~3, p.~341). For every element \(\lambda\) of D and every pair \((i,j)\) of distinct integers in {[}1, n{]}, we set (II, §10, No.~13, p.~361)

\[
B _ {i j} (\lambda) = I _ {n} + \lambda E _ {i j};
\]

it is an element of \(\mathbf{GL}_{n}(\mathrm{D})\) . If \(\lambda_{1},\ldots,\lambda_{n}\) are elements of \(D^{*}\) , then the diagonal matrix \(\mathrm{diag}(\lambda_{1},\ldots,\lambda_{n})\) also belongs to \(\mathbf{GL}_{n}(\mathrm{D})\) .

PROPOSITION 3. --- The mapping det is the unique homomorphism from \(\mathbf{GL}_{n}(\mathrm{D})\) to \(D_{ab}^{*}\) that satisfies the relations

\[
\det (B _ {i j} (1)) = 1\tag{16}
\]

for \(i\neq j\) and

\[
\det (\mathrm{diag} (\lambda_ {1}, \ldots , \lambda_ {n})) = \pi (\lambda_ {1} \dots \lambda_ {n}),\tag{17}
\]

for \(\lambda_1, \ldots, \lambda_n \in \mathrm{D}^*\)

Let \(A\) be an element of \(\mathbf{GL}_n(\mathrm{D})\) and \(a_1, \ldots, a_n\) be its columns. We have

\[
\det (A) = \omega_ {0} (a _ {1}, \dots , a _ {n}).
\]

Now, the columns of the matrix \(A \operatorname{diag}(\lambda_1, \ldots, \lambda_n)\) are \(a_1\lambda_1, \ldots, a_n\lambda_n\) , and those of the matrix \(AB_{ij}(1)\) are therefore \(a_1, \ldots, a_j + a_i, a_{j+1}, \ldots, a_n\) . Since \(\omega_0\) is the unique element of \(\Omega(\mathrm{D}_d^n)\) such that \(\omega_0(\varepsilon_1, \ldots, \varepsilon_n) = 1\) , we see that the determinant is the unique mapping \(\varphi: \mathbf{GL}_n(\mathrm{D}) \to \mathrm{D}_{\mathrm{ab}}^*\) that satisfies the relations

\begin{enumerate}
\def\labelenumi{(\arabic{enumi})}
\setcounter{enumi}{17}
\tightlist
\item
\end{enumerate}

\[
\varphi (A \operatorname{diag} (\lambda_ {1}, \ldots , \lambda_ {n})) = \varphi (A) \pi (\lambda_ {1} \dots \lambda_ {n}),\tag{19}
\]

\[
\varphi (A B _ {i j} (1)) = \varphi (A) \quad \mathrm{for} i \neq j,\tag{20}
\]

\[
\varphi (I _ {n}) = 1.
\]

This proves, to begin with, that the determinant mapping satisfies relations (16) and (17). We already know that this mapping is a homomorphism from \(\mathbf{GL}_{n}(\mathrm{D})\) to \(D_{ab}^{*}\) . Conversely, if \(\varphi\) is a homomorphism from \(\mathbf{GL}_{n}(\mathrm{D})\) to \(D_{ab}^{*}\) such that \(\varphi(B_{ij}(1)) = 1\) for \(i \neq j\) and \(\varphi(\text{diag}(\lambda_1, \ldots, \lambda_n)) = \pi(\lambda_1 \cdots \lambda_n)\) , then \(\varphi\) satisfies relations (18) through (20) and is therefore equal to det.

Examples. --- 1) The columns of the matrix \(B_{ij}(\lambda)\) are \(\varepsilon_1, \ldots, \varepsilon_j + \varepsilon_i \lambda, \varepsilon_{i+1}, \ldots, \varepsilon_n\) . By formula (5) of VIII, p.~448, we have

\[
\det (B _ {i j} (\lambda)) = 1.\tag{21}
\]

\begin{enumerate}
\def\labelenumi{\arabic{enumi})}
\setcounter{enumi}{1}
\tightlist
\item
  Let \(\sigma\) be a permutation of the interval \([1,n]\) in \(\mathbf{N}\) , with signature \(\varepsilon(\sigma)\) . Let \(M(\sigma)\) be the matrix of the permutation \(\sigma\) (II, §10, No.~7, p.~351). The columns of the matrix \(M(\sigma)\) are \(\varepsilon_{\sigma(1)},\ldots,\varepsilon_{\sigma(n)}\) . Applying formula (4) of VIII, p.~448, we therefore have
\end{enumerate}

\[
\det (M (\sigma)) = \pi (\varepsilon (\sigma)).\tag{22}
\]

\begin{enumerate}
\def\labelenumi{\arabic{enumi})}
\setcounter{enumi}{2}
\item
  Suppose \(n \geqslant 1\) . For every invertible diagonal matrix of the form \(\Delta = \text{diag}(d_1, \ldots, d_n)\) , we have \(\Delta B_{ij}(\lambda)\Delta^{-1} = B_{ij}(d_i\lambda d_j^{-1})\) . Let A be an element of \(\mathbf{GL}_n(\mathrm{D})\) . By Corollary 1 of II, §10, No.~13, p.~362 and the previous formula, there exist matrices P and \(\Delta\) in \(\mathbf{GL}_n(\mathrm{D})\) such that \(A = P\Delta\) , that P is the product of matrices of the form \(B_{ij}(\lambda)\) , and that \(\Delta\) is a diagonal matrix of the form \(\text{diag}(1, \ldots, 1, d)\) . We have \(\det(P) = 1\) by Example 1, and therefore \(\det(A) = \det(\Delta) = \pi(d)\) by Proposition 3.
\item
  Let \(D'\) be a field, and let u be a homomorphism from D to \(D'\) . When passing to the quotients, u defines a group homomorphism \(u_{ab}\) from \(D_{ab}^{*}\) to \(D_{ab}^{\prime *}\) . Let \(u_{n}\) be the homomorphism from \(\mathbf{GL}_{n}(\mathrm{D})\) to \(\mathbf{GL}_{n}(\mathrm{D}')\) that transforms a matrix \(A = (a_{ij})\) into the matrix \((u(a_{ij}))\) . The formula
\end{enumerate}

\[
\det (u _ {n} (A)) = u _ {\mathrm{ab}} (\det (A))\tag{23}
\]

for \(A \in \mathbf{GL}_n(\mathrm{D})\) follows immediately from Example 3.

\begin{enumerate}
\def\labelenumi{\arabic{enumi})}
\setcounter{enumi}{4}
\item
  We have \(\mathbf{GL}_1(\mathrm{D}) = \mathrm{D}^*\) , and Proposition 3 shows that we have \(\det(a) = \pi(a)\) for \(a\) in \(\mathbf{GL}_1(\mathrm{D})\) .
\item
  Suppose \(n = 2\) . Let \(A = \begin{pmatrix} a & b \\ c & d \end{pmatrix}\) be an element of \(\mathbf{GL}_2(\mathrm{D})\) . The elements \(a\) and \(c\) of \(\mathrm{D}\) are not both zero. We will give the determinant of \(A\) explicitly.
\end{enumerate}

If \(a\) is not zero, then we have

\[
A = \left( \begin{array}{c c} 1 & 0 \\ c a ^ {- 1} & 1 \end{array} \right) \left( \begin{array}{c c} a & 0 \\ 0 & d - c a ^ {- 1} b \end{array} \right) \left( \begin{array}{c c} 1 & a ^ {- 1} b \\ 0 & 1 \end{array} \right).\tag{24}
\]

Consequently, \(ad - aca^{-1}b \neq 0\) and

\[
\det (A) = \pi (a d - a c a ^ {- 1} b).\tag{25}
\]

\begin{enumerate}
\def\labelenumi{\alph{enumi})}
\setcounter{enumi}{1}
\tightlist
\item
  If \(a\) is zero, then we have \(c \neq 0\) and
\end{enumerate}

\[
A = \left( \begin{array}{c c} 0 & b \\ c & d \end{array} \right) = \left( \begin{array}{c c} 0 & 1 \\ 1 & 0 \end{array} \right) \left( \begin{array}{c c} c & d \\ 0 & b \end{array} \right).\tag{26}
\]

Consequently, by a) and Example 2, we have \(cb \neq 0\) and

\[
\det (A) = \pi (- c b).\tag{27}
\]

\begin{enumerate}
\def\labelenumi{\arabic{enumi})}
\setcounter{enumi}{6}
\tightlist
\item
  Let \(D^{o}\) be the opposite field of D. The mapping \(A \mapsto {}^{t}A\) from \(\mathbf{M}_{n}(\mathrm{D})\) to \(\mathbf{M}_{n}(\mathrm{D}^{\circ})\) satisfies \({}^{t}(AB) = {}^{t}B^{t}A\) . Therefore, for every matrix A in \(\mathbf{GL}_{n}(\mathrm{D})\) , the transposed matrix \({}^{t}A\) is invertible in \(\mathbf{M}_{n}(\mathrm{D}^{\circ})\) , but it is not necessarily invertible in \(\mathbf{M}_{n}(\mathrm{D})\) . It follows from Example 3 that if A belongs to \(\mathbf{GL}_{n}(\mathrm{D})\) , then the elements A of \(\mathbf{GL}_{n}(\mathrm{D})\) and \({}^{t}A\) of \(\mathbf{GL}_{n}(\mathrm{D}^{\circ})\) have the same determinant. On the other hand, even if the matrix \({}^{t}A\) belongs to \(\mathbf{GL}_{n}(\mathrm{D})\) , its determinant in \(\mathrm{GL}_{n}(\mathrm{D})\) is not necessarily equal to that of A (cf.~Example 6).
\end{enumerate}

Remarks. --- 1) Let V be a right vector space over the field D, of finite dimension n.~We view the dual space \(V^{*} = \operatorname{Hom}_{\mathrm{D}}(\mathrm{V}, \mathrm{D})\) as a right vector space over the opposite field \(D^{o}\) of D. Let u be an automorphism of V, and let \(^{t}u\) be the automorphism of \(V^{*}\) that is the transpose of u. If u is represented by a matrix A in \(\mathbf{M}_{n}(\mathrm{D})\) with respect to a basis \((e_{1}, \ldots, e_{n})\) of V, then the automorphism \(^{t}u\) is represented by the matrix \(^{t}A\) in \(\mathbf{M}_{n}(\mathrm{D}^{o})\) with respect to the basis \((e_{1}^{*}, \ldots, e_{n}^{*})\) of \(V^{*}\) dual to \((e_{1}, \ldots, e_{n})\) (II, §10, No.~4, p.~344, Proposition 3). By Example 7, the determinant of \(^{t}u\) is equal to that of u.

\begin{enumerate}
\def\labelenumi{\arabic{enumi})}
\setcounter{enumi}{1}
\tightlist
\item
  The results of Subsections 1 through 4 generalize to the case when D is a local ring (VIII, p.~459, Exercise 2).
\end{enumerate}

\section{The Unimodular Group}

Let n be a positive integer. We denote by \(\mathbf{SL}_{n}(\mathrm{D})\) the kernel of the group homomorphism \(\det: \mathbf{GL}_{n}(\mathrm{D}) \to \mathrm{D}_{\mathrm{ab}}^{*}\) , and we call it the unimodular group or special linear group (cf.~III, §8, No.~9, p.~537 when the field D is commutative).

THEOREM 1. --- Suppose \(n \geqslant 2\) .

\begin{enumerate}
\def\labelenumi{\alph{enumi})}
\item
  The subgroup \(\mathbf{SL}_n(\mathrm{D})\) of \(\mathbf{GL}_n(\mathrm{D})\) is generated by the matrices \(B_{ij}(\lambda)\) , where \(i\) and \(j\) are distinct integers in the interval \([1,n]\) and \(\lambda\) runs through D.
\item
  Suppose \(n \geqslant 3\) or \(\operatorname{Card}(\mathrm{D}) \geqslant 3\) . The derived group of \(\mathbf{GL}_n(\mathrm{D})\) is equal to \(\mathbf{SL}_n(\mathrm{D})\) .
\item
  Suppose \(n \geqslant 3\) or \(\operatorname{Card}(\mathrm{D}) \geqslant 4\) . The derived group of \(\mathbf{SL}_n(\mathrm{D})\) is equal to \(\mathbf{SL}_n(\mathrm{D})\) .
\end{enumerate}

\begin{enumerate}
\def\labelenumi{\Alph{enumi})}
\tightlist
\item
  Denote by T the subgroup of \(\mathbf{GL}_{n}(\mathrm{D})\) generated by the matrices \(B_{ij}(\lambda)\) . By Example 1 of VIII, p.~454, we have \(\det(B_{ij}(\lambda)) = 1\) , and therefore \(\mathrm{T} \subset \mathbf{SL}_{n}(\mathrm{D})\) . To prove that these two groups are equal, it then suffices, by Example 3 of loc. cit., to prove that every matrix of the form \(\operatorname{diag}(1, \ldots, 1, d)\) with \(\pi(d) = 1\) belongs to T. The matrix \(\operatorname{diag}(1, \ldots, 1, d)\) belongs to the image of the homomorphism \(U \mapsto \begin{pmatrix} I_{n-2} & 0 \\ 0 & U \end{pmatrix}\) from \(\mathbf{GL}_{2}(\mathrm{D})\) to \(\mathbf{GL}_{n}(\mathrm{D})\) ; it therefore suffices to consider the case n = 2. Since the kernel of \(\pi\) is the derived group of \(D^{*}\) , we may assume \(d = uvu^{-1}v^{-1}\) with u, v in \(D^{*}\) . Our assertion then follows from the equalities
\end{enumerate}

\[
\left( \begin{array}{c c} 1 & 0 \\ 0 & d \end{array} \right) = \left( \begin{array}{c c} u ^ {- 1} & 0 \\ 0 & u \end{array} \right) \left( \begin{array}{c c} v ^ {- 1} & 0 \\ 0 & v \end{array} \right) \left( \begin{array}{c c} v u & 0 \\ 0 & u ^ {- 1} v ^ {- 1} \end{array} \right)\tag{28}
\]

and

\[
\left( \begin{array}{c c} s & 0 \\ 0 & s ^ {- 1} \end{array} \right) = \left( \begin{array}{c c} 1 & s \\ 0 & 1 \end{array} \right) \left( \begin{array}{c c} 1 & 0 \\ 1 - s ^ {- 1} & 1 \end{array} \right) \left( \begin{array}{c c} 1 & - 1 \\ 0 & 1 \end{array} \right) \left( \begin{array}{c c} 1 & 0 \\ 1 - s & 1 \end{array} \right)\tag{29}
\]

for \(s \in D^{*}\) .

\begin{enumerate}
\def\labelenumi{\Alph{enumi})}
\setcounter{enumi}{1}
\tightlist
\item
  By construction, \(\mathbf{SL}_{n}(\mathrm{D})\) is the kernel of a homomorphism from \(\mathbf{GL}_{n}(\mathrm{D})\) to a commutative group, so it contains the derived group \((\mathbf{GL}_{n}(\mathrm{D}),\mathbf{GL}_{n}(\mathrm{D}))\) of \(\mathbf{GL}_{n}(\mathrm{D})\) . By the above, we have
\end{enumerate}

\[
\mathbf {S L} _ {n} (\mathrm{D}) \supset (\mathbf {G L} _ {n} (\mathrm{D}), \mathbf {G L} _ {n} (\mathrm{D})) \supset (\mathbf {S L} _ {n} (\mathrm{D}), \mathbf {S L} _ {n} (\mathrm{D})).
\]

To prove c), it suffices to prove that the matrices \(B_{ij}(\lambda)\) are commutators in \(\mathbf{SL}_n(\mathrm{D})\) .

Suppose \(n \geqslant 3\) . If \(i, j, k\) are distinct integers in the interval \([1, n]\) and \(\mu, \nu\) are elements of \(D\) , then we have

\[
B _ {i j} (\mu \nu) = B _ {i k} (\mu) ^ {- 1} B _ {k j} (\nu) ^ {- 1} B _ {i k} (\mu) B _ {k j} (\nu).\tag{30}
\]

By taking \(\mu = 1\) and \(\nu = \lambda\) , we see that the matrix \(B_{ij}(\lambda)\) is a commutator of elements of \(\mathbf{SL}_n(\mathrm{D})\) .

Now suppose \(n = 2\) . Let \(u\) and \(v\) be elements of \(D\) with \(u \neq 0\) . We have the relations

\[
\left( \begin{array}{c c} 1 & v - u v u \\ 0 & 1 \end{array} \right) = \left( \begin{array}{c c} u & 0 \\ 0 & u ^ {- 1} \end{array} \right) \left( \begin{array}{c c} 1 & - v \\ 0 & 1 \end{array} \right) \left( \begin{array}{c c} u ^ {- 1} & 0 \\ 0 & u \end{array} \right) \left( \begin{array}{c c} 1 & v \\ 0 & 1 \end{array} \right)\tag{31}
\]

and

\[
\left( \begin{array}{c c} 1 & 0 \\ v - u v u & 1 \end{array} \right) = \left( \begin{array}{c c} u ^ {- 1} & 0 \\ 0 & u \end{array} \right) \left( \begin{array}{c c} 1 & 0 \\ - v & 1 \end{array} \right) \left( \begin{array}{c c} u & 0 \\ 0 & u ^ {- 1} \end{array} \right) \left( \begin{array}{c c} 1 & 0 \\ v & 1 \end{array} \right).\tag{32}
\]

We have \(\det\left(\begin{matrix}u & 0 \\ 0 & u^{-1}\end{matrix}\right)=1\) , so the matrices \(B_{12}(v-uvu)\) and \(B_{21}(v-uvu)\) are commutators of elements of \(\mathbf{SL}_{n}(\mathrm{D})\) .

Suppose that the field D has at least four elements. Let \(\lambda\) be an element of D. If \(\lambda\) is equal to 0, 1, or -1, choose an arbitrary element u in \(D-\{0,1,-1\}\) ; otherwise, set \(u=\lambda\) . In both cases, u is a nonzero element of D, it commutes with \(\lambda\) , and we have \(u^{2}\neq1\) . Set \(v=\lambda(1-u^{2})^{-1}\) . We have uv=vu and therefore \(v-uvu=v(1-u^{2})=\lambda\) . It then follows from relations (31) and (32) that the matrices \(B_{12}(\lambda)\) and \(B_{21}(\lambda)\) are commutators in \(\mathbf{SL}_{n}(\mathrm{D})\) . This completes the proof of c).

\begin{enumerate}
\def\labelenumi{\Alph{enumi})}
\setcounter{enumi}{2}
\tightlist
\item
  It remains to prove that \(\mathbf{SL}_{2}(\mathrm{D})\) is the derived group of \(\mathbf{GL}_{2}(\mathrm{D})\) when D has three elements. By A), the group \(\mathbf{SL}_{2}(\mathrm{D})\) is generated by the matrices \(B_{12}(1)=(\begin{matrix}{1}&{1}\\ {0}&{1}\end{matrix})\) and \(B_{21}(1)=(\begin{matrix}{1}&{0}\\ {1}&{1}\end{matrix})\) , and these matrices are commutators of elements of \(\mathbf{GL}_{2}(\mathrm{D})\) because we have \(B_{21}(1)={}^{t}B_{12}(1)\) and
\end{enumerate}

\[
\left( \begin{array}{c c} 1 & 1 \\ 0 & 1 \end{array} \right) = \left( \begin{array}{c c} - 1 & 0 \\ 0 & 1 \end{array} \right) \left( \begin{array}{c c} 1 & 1 \\ 0 & 1 \end{array} \right) \left( \begin{array}{c c} - 1 & 0 \\ 0 & 1 \end{array} \right) ^ {- 1} \left( \begin{array}{c c} 1 & 1 \\ 0 & 1 \end{array} \right) ^ {- 1}.\tag{33}
\]

Remark. --- If D is a field with two elements, then \(\mathbf{GL}_{2}(\mathrm{D})\) is equal to \(\mathbf{SL}_{2}(\mathrm{D})\) , and it is a group of order 6 whose derived group has order 3. If D is a field with three elements, then the group \(\mathbf{SL}_{2}(\mathrm{D})\) has order 24, and its derived group has order 8. Suppose \(n \geqslant 2\) ; except in the two previous cases, every normal subgroup of \(\mathbf{SL}_{n}(\mathrm{D})\) distinct from \(\mathbf{SL}_{n}(\mathrm{D})\) is contained in the center \(\mathrm{Z}(\mathbf{SL}_{n}(\mathrm{D}))\) of \(\mathbf{SL}_{n}(\mathrm{D})\) (II, §10, p.~421, Exercises 13 and 14). The group \(\mathbf{SL}_{n}(\mathrm{D})/\mathrm{Z}(\mathbf{SL}_{n}(\mathrm{D}))\) is then simple.

Let V be a right vector space over the field D, of finite dimension n.~We denote by \(\mathrm{SL}(V)\) , and call unimodular group of V, the kernel of the homomorphism \(\det: \mathbf{GL}(V) \to D_{ab}^{*}\) . Choosing a basis of V allows one to identify this group with \(\mathbf{SL}_{n}(D)\) .

We say that an automorphism u of V is a transvection if there exist a vector v of V and a linear form \(\varphi\) on V such that \(\varphi(v)=0\) and \(u(x)=x+v\varphi(x)\) for every x in V. When \(n\geqslant2\) , u is a transvection if and only if there exists a basis of V with respect to which the matrix of u is of the form \(B_{ij}(\lambda)\) . Theorem 1 implies the following corollary.

COROLLARY. --- Let V be a right vector space over the field D, of finite dimension \(n \geqslant 2\) .

\begin{enumerate}
\def\labelenumi{\alph{enumi})}
\item
  The subgroup \(\mathbf{SL}(\mathrm{V})\) of \(\mathbf{GL}(\mathrm{V})\) is generated by the transvections.
\item
  The subgroup \(\mathbf{SL}(\mathrm{V})\) is the derived group of \(\mathbf{GL}(\mathrm{V})\) unless we have \(n = 2\) and D has two elements.
\item
  The group \(\mathbf{SL}(V)\) is equal to its derived group unless we have n = 2 and D has two or three elements.
\end{enumerate}

\BourbakiAppendixExercises{Exercises}

\begin{enumerate}
\def\labelenumi{\arabic{enumi})}
\item
  Give an example of an invertible matrix A over a field D whose transpose is not invertible, and of an invertible matrix B over D whose transpose is invertible but satisfies \(\det^{t}B \neq \det B\) (consider square matrices of order 2 over a quaternion field).
\item
  Let D be a local ring; we denote by \(D^{*}\) the group of invertible elements of D, by \(D_{ab}^{*}\) the quotient of \(D^{*}\) by its derived group, and by \(\pi : D^{*} \to D_{ab}^{*}\) the canonical homomorphism. Let V be a free right D-module of finite dimension n.~We denote by \(\mathrm{B(V)}\) the set of bases of V and by \(\Omega(\mathrm{V})\) the set of mappings \(\omega : \mathrm{B(V)} \to \mathrm{D}_{\mathrm{ab}}^{*}\) satisfying the following two conditions:
\end{enumerate}

\begin{enumerate}
\def\labelenumi{(\roman{enumi})}
\item
  For \(\lambda_1, \ldots, \lambda_n\) in D* and \((v_1, \ldots, v_n) \in \mathrm{B}(\mathrm{V})\) , we have \(\omega(v_1 \lambda_1, \ldots, v_n \lambda_n) = \omega(v_1, \ldots, v_n)\pi(\lambda_1 \ldots \lambda_n)\) .
\item
  For \(i \neq j\) and \(\lambda \in D\) , we have \(\omega(v_{1},\ldots,v_{i}+v_{j}\lambda,\ldots,v_{n}) = \omega(v_{1},\ldots,v_{n})\) . a) Let e be an element of V and W a submodule of V supplementary to eD. Let \(\varphi \in \Omega(\mathrm{W})\) ; prove that there exists a unique element \(\omega\) of \(\Omega(\mathrm{V})\) satisfying
\end{enumerate}

\[
\omega (w _ {1}, \dots , w _ {n - 1}, e) = \varphi (w _ {1}, \dots , w _ {n - 1})
\]

for every basis \((w_{1},\ldots ,w_{n - 1})\) of W (adapt the proof of Proposition 1).

\begin{enumerate}
\def\labelenumi{\alph{enumi})}
\setcounter{enumi}{1}
\item
  Deduce that given two elements \(\omega\) and \(\omega'\) of \(\Omega(\mathrm{V})\) , there exists a unique element t of \(D_{ab}^{*}\) such that we have \(\omega' = t\omega\) .
\item
  Let u be an automorphism of V. Prove that there exists a unique element of \(D_{ab}^{*}\) , denoted by \(\det u\) , such that we have \(\omega(u(v_{1}),\ldots,u(v_{n}))=(\det u)\omega(v_{1},\ldots,v_{n})\) for every basis \((v_{1},\ldots,v_{n})\) of V and every element \(\omega\) of \(\Omega(\mathrm{V})\) . The mapping \(u\mapsto\det u\) from \(\mathbf{GL}(\mathrm{V})\) to \(D_{ab}^{*}\) is a group homomorphism.
\end{enumerate}

\(d\) ) Take \(\mathrm{V} = \mathrm{D}_d^n\) , so that \(\mathbf{GL}(\mathrm{V})\) is identified with the group \(\mathbf{GL}_n(\mathrm{D})\) of invertible square matrices of order \(n\) . Prove that the mapping \(\operatorname{det}:\mathbf{GL}_n(\mathrm{D})\to \mathrm{D}_{\mathrm{ab}}^*\) is the unique group homomorphism with value 1 on all matrices \(B_{ij}(\lambda)\) ( \(i\neq j\) , \(\lambda \in \mathrm{D}\) ) and satisfying \(\operatorname{det}(\operatorname{diag}(\lambda_1,\dots ,\lambda_n)) = \pi (\lambda_1\dots \lambda_n)\) for \(\lambda_1,\ldots ,\lambda_n\) in \(\mathrm{D}^*\) .

\begin{enumerate}
\def\labelenumi{\alph{enumi})}
\setcounter{enumi}{4}
\tightlist
\item
  Adapt Examples 2 through 6 of No.~4 to this setting.
\end{enumerate}

¶ 3) Let \(\mathrm{A} = \mathrm{A}_0 \oplus \mathrm{A}_1\) be a graded ring of type \(\mathbf{Z}/2\mathbf{Z}\) . Suppose that \(\mathrm{A}_0\) is contained in the center of \(\mathrm{A}\) and that every element of \(\mathrm{A}_1\) has square zero.

\begin{enumerate}
\def\labelenumi{\alph{enumi})}
\item
  Let \(A_{1}^{2}\) be the Z-submodule of \(A_{0}\) generated by the products of two elements of \(A_{1}\) ; show that \(A_{1}^{2}\) is a nil ideal of \(A_{0}\) and that \(A_{1}^{2} \oplus A_{1}\) is a two-sided nil ideal of A.
\item
  Let \(p, q\) be integers \(\geqslant 0\) . Denote by \(A^{p|q}\) the free graded right A-module \(A_d^p \oplus A_d(1)^q\) (II, §11, No.~2, p.~365, Example 3) and by \(\mathbf{GL}(p|q, \mathrm{A})\) the group of (graded) automorphisms (of degree 0) of \(A^{p|q}\) . Such an automorphism is represented in the canonical basis of \(A^{p|q}\) by a matrix \(X = \begin{pmatrix} A & B \\ C & D \end{pmatrix}\) with \(A \in \mathbf{M}_p(\mathrm{A}_0)\) , \(B \in \mathbf{M}_{p,q}(\mathrm{A}_1)\) , \(C \in \mathbf{M}_{q,p}(\mathrm{A}_1)\) , \(D \in \mathbf{M}_q(\mathrm{A}_0)\) . Prove that the matrices \(A\) and \(D\) are invertible (use \(a\) ) and Exercise 16 of VIII, p.~168).
\item
  For X as above, set \(\operatorname{sdet} X = \det(A - BD^{-1}C) \det D^{-1}\) . Prove that \(\operatorname{sdet} X\) is an invertible element of \(A_{0}\) .
\end{enumerate}

\(d\) ) Prove that every matrix \(X\in \mathbf{GL}(p|q,\mathrm{A})\) can be written as a product

\[
X = \left( \begin{array}{c c} I & R \\ 0 & I \end{array} \right) \left( \begin{array}{c c} P & 0 \\ 0 & Q \end{array} \right) \left( \begin{array}{c c} I & 0 \\ S & I \end{array} \right).
\]

(Take \(R = BD^{-1}\) , \(P = A - BD^{-1}C\) , \(Q = D\) , \(S = D^{-1}C\) .)

\begin{enumerate}
\def\labelenumi{\alph{enumi})}
\setcounter{enumi}{4}
\tightlist
\item
  Prove that sdet is a homomorphism from \(\mathbf{GL}(p|q,\mathrm{A})\) to the group of invertible elements of \(A_{0}\) . (Use d) to reduce to proving the equality \(\operatorname{sdet}(XY)=\operatorname{sdet}X\operatorname{sdet}Y\) when X is of the form \(\binom{I}{S}\operatorname{sdet}(I)\) , where S is an elementary matrix. Observe that for every matrix \(R\in\mathbf{M}_{p,q}(\mathrm{A}_{1})\) , the product of two arbitrary entries of the matrix RS (resp. SR) is zero, and deduce the equality \(\det(I-RS)=(\det(I-SR))^{-1}\) .
\item
  Let \(X = \begin{pmatrix} A & B \\ C & D \end{pmatrix}\) be an element of \(\mathbf{GL}(p|q,\mathrm{A})\) ; set \(^{st}X = \begin{pmatrix} t_{A} & t_{C} \\ -t_{B} & t_{D} \end{pmatrix}\) and \(^{\pi}X = \begin{pmatrix} D & C \\ B & A \end{pmatrix}\) . Let \(Y \in \mathbf{GL}(p|q, A)\) ; we have \(^{st}(XY) = ^{st}Y^{st}X\) and \(^{\pi}(XY) = ^{\pi}X^{\pi}Y\) . Prove the equalities \(\operatorname{sdet}X = \operatorname{sdet}^{st}X = (\operatorname{sdet}^{\pi}X)^{-1}\) (use d) to calculate \(\operatorname{sdet}^{\pi}X\) .
\end{enumerate}

\begin{enumerate}
\def\labelenumi{\arabic{enumi})}
\setcounter{enumi}{3}
\tightlist
\item
  Keep the assumptions of Exercise 3. Let \(\mathrm{M} = \mathrm{M}_0 \oplus \mathrm{M}_1\) be a finitely generated free graded right A-module; there exist unique integers \(p, q\) such that \(\mathrm{M}\) is isomorphic to \(\mathrm{A}^{p|q}\) . We say that a basis of \(\mathrm{M}\) is adapted if its first \(p\) prime vectors are homogeneous of degree 0 and the last \(q\) are homogeneous of degree 1.
\end{enumerate}

\begin{enumerate}
\def\labelenumi{\alph{enumi})}
\item
  Let u be an automorphism of the graded module M, and let \(M_{\mathcal{B}}(u)\) be its matrix with respect to an adapted basis of B. Prove that the element sdet \(M_{\mathcal{B}}(u)\) does not depend on the choice of B. We denote it by sdet u.
\item
  Denote by \(u(1)\) the endomorphism u viewed as an automorphism of the graded module M(1). If B is an adapted basis of M, then the basis \(^{\pi}B\) obtained by interchanging the first p vectors of B and the last q is an adapted basis of M(1), and we have \(M_{\pi\mathcal{B}}(u(1)) = ^{\pi}M_{\mathcal{B}}(u)\) (Exercise 3, f)). Deduce that we have \(\operatorname{sdet}(u(1)) = (\operatorname{sdet} u)^{-1}\) .
\item
  On every graded left A-module N, we define a canonical graded right A-module structure by setting \(ax = (-1)^{ij}xa\) for \(a \in A_i\) , \(x \in N_j\) . In particular, we view the graded dual \(M^{*gr}\) of M as a graded right A-module, also isomorphic to \(A^{p|q}\) . Let u be an automorphism of M. We denote by \({}^{st}u\) the automorphism of \(M^{*gr}\) characterized by the formula \(\langle{}^{st}u(x^{*}), x\rangle = (-1)^{rs}\langle x^{*}, u(x)\rangle\) for \(x^{*} \in (\mathrm{M}^{*\mathrm{gr}})_r\) , \(x \in M_s\) . If B is an adapted basis of M, then the dual basis \(B^*\) is an adapted basis of \(M^{*gr}\) , and we have \(M_{\mathcal{B}^*}({}^{st}u) = {}^{st}M_{\mathcal{B}}(u)\) (Exercise 3, f). We have \(\operatorname{sdet}({}^{st}u) = \operatorname{sdet}(u)\) .
\end{enumerate}

\chapter{Hilbert's Nullstellensatz}

THEOREM 1. --- Let A be an integral domain, K be its field of fractions, and L be a commutative K-algebra. Suppose that the A-algebra L is generated by finitely many elements and that L is a field.

\begin{enumerate}
\def\labelenumi{\alph{enumi})}
\item
  The degree of L over K is finite.
\item
  There exists a nonzero element \(a\) of \(\mathrm{A}\) such that \(\mathrm{K}\) is equal to \(\mathrm{A}[a^{-1}]\) .
\end{enumerate}

Let \(S\) be a generating subset of the A-algebra \(L\) . We use induction on the cardinal of \(S\) .

First, suppose that the elements of S are all algebraic over K. The degree of L over K is then finite (V, §3, No.~2, p.~18, Theorem 2). Let \((e_i)_{i\in I}\) be a basis of L over K. There exists a nonzero element \(a\) of A such that the coordinates with respect to this basis of the elements of S, of the element 1, and of the elements \(e_i e_j\) for \(i\) , \(j\) in I belong to \(\mathrm{A}[a^{-1}]\) . The set of linear combinations \(\sum_{i\in I}a_i e_i\) with \(a_i\in \mathrm{A}[a^{-1}]\) for every \(i\) is then a subring of L. It contains A and S by construction and is therefore equal to L. In particular, it contains \(\mathrm{Ke}_1\) , so K is equal to \(\mathrm{A}[a^{-1}]\) .

Now suppose that an element \(s\) of \(S\) is transcendental over \(K\) . Denote by \(E\) the field of fractions of the ring \(A[s]\) . The \(A[s]\) -algebra \(L\) is generated by \(S - \{s\}\) . By the induction hypothesis, there exists a nonzero polynomial \(P \in A[X]\) such that \(E\) is equal to \(A[s][P(s)^{-1}]\) . Let \(\overline{K}\) be an algebraic closure of \(K\) (V, §4, No.~3, p.~23, Theorem 2). Since the field \(\overline{K}\) is infinite (V, §4, No.~1, p.~20, Proposition 3), there exists an element \(x\) of \(\overline{K}\) such that \(P(x) \neq 0\) . Let \(\varphi : E \to \overline{K}\) be the unique homomorphism from the K-algebra \(E = A[s][P(s)^{-1}]\) to the K-algebra \(\overline{K}\) that sends \(s\) to \(x\) . This is absurd because \(E\) is a transcendental extension of \(K\) and \(\overline{K}\) is an algebraic extension of \(K\) . This completes the proof of the theorem.

COROLLARY 1. --- Let K be a commutative field, A be a commutative K-algebra generated by finitely many elements, and m be a maximal ideal of A. a) The degree of A/m over K is finite.

\begin{enumerate}
\def\labelenumi{\alph{enumi})}
\setcounter{enumi}{1}
\tightlist
\item
  Let \(\Omega\) be an algebraically closed extension of K. There exists a K-algebra homomorphism from A to \(\Omega\) with kernel \(\mathfrak{m}\) .
\end{enumerate}

The K-algebra A/m is generated by finitely many elements, and A/m is a field. By Theorem 1, the degree of A/m is finite; assertion a) follows. Every extension of K of finite degree is isomorphic to a subextension of \(\Omega\) (V, §4, No.~1, p.~20, Theorem 1); this gives b).

COROLLARY 2. --- Let K be a commutative field, n be a natural number, \((\mathrm{P}_{i})_{i\in\mathrm{I}}\) be a family of elements of \(K[X_{1},\ldots,X_{n}]\) , and \(\Omega\) be an algebraically closed extension of K. The following conditions are equivalent:

\begin{enumerate}
\def\labelenumi{(\roman{enumi})}
\item
  The polynomials \(\mathrm{P}_i\) do not have a common zero in \(\Omega^n\) .
\item
  There exists a family \((\mathrm{Q}_i)_{i\in \mathrm{I}}\) , with finite support, of elements of \(\mathrm{K}[\mathrm{X}_1,\dots ,\mathrm{X}_n]\) such that \(\sum_{i\in \mathrm{I}}\mathrm{P}_i\mathrm{Q}_i = 1\) .
\end{enumerate}

Denote by A the ring \(K[X_{1},\ldots,X_{n}]\) and by a the ideal generated by the polynomial \(P_{i}\) . Condition (i) means that there are no K-algebra homomorphisms from A/a to \(\Omega\) . If it is satisfied, then by Corollary 1, the ring A/a has no maximal ideals and is therefore zero, and 1 belongs to a. This proves that (i) implies (ii). The implication (ii) \(\Rightarrow\) (i) is clear.

\chapter{Trace of an Endomorphism of Finite Rank}

\section{Linear Mappings of Finite Rank}

Let A be a ring, and let E and F be A-modules. We denote by \(\operatorname{Hom}_{\mathrm{A}}^{\mathrm{f}}(\mathrm{E},\mathrm{F})\) the set of linear mappings from E to F whose image is contained in a finitely generated submodule of F. It is a subgroup of \(\operatorname{Hom}_{\mathrm{A}}(\mathrm{E},\mathrm{F})\) . When A is a field, it is the set of linear mappings from E to F of finite rank (II, §7, No.~4, p.~298, Definition 3). Set \(\operatorname{End}_{\mathrm{A}}^{\mathrm{f}}(\mathrm{E}) = \operatorname{Hom}_{\mathrm{A}}^{\mathrm{f}}(\mathrm{E},\mathrm{E})\) .

Let E, F, G be A-modules, and let \(u : E \to F\) and \(v : F \to G\) be A-linear mappings. If \(u \in \operatorname{Hom}_{\mathrm{A}}^{\mathrm{f}}(\mathrm{E}, \mathrm{F})\) or \(v \in \operatorname{Hom}_{\mathrm{A}}^{\mathrm{f}}(\mathrm{F}, \mathrm{G})\) , then \(v \circ u\) belongs to \(\operatorname{Hom}_{\mathrm{A}}^{\mathrm{f}}(\mathrm{E}, \mathrm{G})\) .

Denote by \(\theta\) the canonical group homomorphism from \(\mathrm{E}^* \otimes_{\mathrm{A}} \mathrm{F}\) to \(\mathrm{Hom}_{\mathrm{A}}(\mathrm{E}, \mathrm{F})\) (II, §4, No.~2, p.~271); it sends an element \(x^* \otimes y\) of \(\mathrm{E}^* \otimes_{\mathrm{A}} \mathrm{F}\) to the mapping \(x \mapsto \langle x, x^* \rangle y\) from \(\mathrm{E}\) to \(\mathrm{F}\) .

Lemma 1. --- Suppose that the A-module F is projective. The group homomorphism \(\theta\) is injective, and its image is \(\mathrm{Hom}_{\mathrm{A}}^{\mathrm{f}}(\mathrm{E},\mathrm{F})\) .

By loc. cit., Corollary, the homomorphism \(\theta\) is injective. Its image is contained in \(\operatorname{Hom}_{\mathrm{A}}^{\mathrm{f}}(\mathrm{E},\mathrm{F})\) . Let \(u \in \operatorname{Hom}_{\mathrm{A}}^{\mathrm{f}}(\mathrm{E},\mathrm{F})\) ; let us prove that u belongs to the image of \(\theta\) .

First, suppose that the A-module F is free. Let \((f_{i})_{i\in\mathrm{I}}\) be a basis of F, and let \((f_{i}^{*})_{i\in\mathrm{I}}\) be the basis of \(F^{*}\) dual to \((f_{i})_{i\in\mathrm{I}}\) . Let J be a finite subset of I such that the image of u is contained in the submodule of F generated by the \(f_{j}\) for \(j\in J\) . We have

\[
u (x) = \sum_ {j \in \mathrm{J}} \langle u (x), f _ {j} ^ {*} \rangle f _ {j} = \sum_ {j \in \mathrm{J}} \langle x, ^ {t} u (f _ {j} ^ {*}) \rangle f _ {j}\tag{1}
\]

for every \(x\in \mathrm{E}\) , and therefore

\[
u = \theta \left(\sum_ {j \in \mathrm{J}} ^ {t} u \left(f _ {j} ^ {*}\right) \otimes f _ {j}\right).\tag{2}
\]

In the general case, there exist a free A-module L and homomorphisms \(i: F \to L\) and \(p: L \to F\) such that \(p \circ i = 1_{F}\) . The homomorphism \(i \circ u\) belongs to \(\operatorname{Hom}_{\mathrm{A}}^{\mathrm{f}}(\mathrm{E}, \mathrm{L})\) . By the above, there exist a finite set J and, for every \(j \in J\) , elements \(x_{j}^{*}\) of \(E^{*}\) and \(y_{j}\) of L such that we have

\[
i (u (x)) = \sum_ {j \in \mathrm{J}} \langle x, x _ {j} ^ {*} \rangle y _ {j}
\]

for every \(x \in \mathrm{E}\) . We then have

\[
u (x) = p \left(i (u (x))\right) = \sum_ {j \in J} \langle x, x _ {j} ^ {*} \rangle p \left(y _ {j}\right)
\]

for every \(x \in \mathrm{E}\) , and therefore \(u = \theta\big(\sum_{j \in \mathrm{J}} x_j^* \otimes p(y_j)\big)\) .

\section{Trace of an Endomorphism of Finite Rank}

In this subsection, A denotes a commutative ring. Let E be a projective A-module. Then \(\operatorname{End}_{\mathrm{A}}^{\mathrm{f}}(\mathrm{E})\) is an A-submodule of \(\operatorname{End}_{\mathrm{A}}(\mathrm{E})\) , and the canonical mapping \(E^{*}\otimes_{A}E\to\operatorname{End}_{\mathrm{A}}(\mathrm{E})\) defines an isomorphism \(\theta_{E}\) of A-modules from \(E^{*}\otimes_{A}E\) to \(\operatorname{End}_{\mathrm{A}}^{\mathrm{f}}(\mathrm{E})\) (Lemma 1). Consider the canonical linear form \(\tau:E^{*}\otimes_{A}E\to A\) (II, §4, No.~3, p.~273) characterized by the formula \(\tau(x^{*}\otimes x)=\langle x,x^{*}\rangle\) . By composing it with the isomorphism \(\theta_{E}^{-1}\) , we deduce a linear form \(\operatorname{Tr}: \operatorname{End}_{\mathrm{A}}^{\mathrm{f}}(\mathrm{E})\to\mathrm{A}\) , called the trace form. When the A-module E is finitely generated, we recover the definition of II, §4, No.~3, p.~273.

PROPOSITION 1. --- Let E be a free A-module. Let \((e_{i})_{i\in\mathrm{I}}\) be a basis of E, and let \((e_{i}^{*})_{i\in\mathrm{I}}\) be its dual basis. Let \(u\in\operatorname{End}_{\mathrm{A}}^{\mathrm{f}}(\mathrm{E})\) . The family \((\langle u(e_{i}),e_{i}^{*}\rangle)_{i\in\mathrm{I}}\) has finite support, and its sum is equal to \(\operatorname{Tr}(u)\) .

It suffices to treat the case when \(u\) is of the form \(\theta_{\mathrm{E}}(x^{*} \otimes x)\) with \(x \in \mathrm{E}\) and \(x^{*} \in \mathrm{E}^{*}\) . The family \((\langle x, e_{i}^{*} \rangle)_{i \in \mathrm{I}}\) then has finite support, and we have \(x = \sum_{i \in \mathrm{I}} \langle x, e_{i}^{*} \rangle e_{i}\) . Consequently, the family \((\langle x, e_{i}^{*} \rangle \langle e_{i}, x^{*} \rangle)_{i \in \mathrm{I}}\) also has finite support, and we have \(\langle x, x^{*} \rangle = \sum_{i \in \mathrm{I}} \langle x, e_{i}^{*} \rangle \langle e_{i}, x^{*} \rangle\) . Now, we have \(\langle u(e_{i}), e_{i}^{*} \rangle = \langle x, e_{i}^{*} \rangle \langle e_{i}, x^{*} \rangle\) for every \(i \in \mathrm{I}\) . This proves the proposition.

PROPOSITION 2. --- Let E, F be projective A-modules. Let \(u \in \operatorname{Hom}_{\mathrm{A}}^{\mathrm{f}}(\mathrm{E}, \mathrm{F})\) and \(v \in \operatorname{Hom}_{\mathrm{A}}(\mathrm{F}, \mathrm{E})\) . We have the relation

\[
\operatorname{Tr} (v \circ u) = \operatorname{Tr} (u \circ v).\tag{3}
\]

It suffices to prove the proposition when \(u\) is of the form \(\theta(x^{*} \otimes y)\) with \(x^{*} \in \mathrm{E}^{*}\) and \(y \in \mathrm{F}\) . In this case, we have

\[
v \circ u = \theta_ {\mathrm{E}} (x ^ {*} \otimes v (y)) \quad \text { and } \quad u \circ v = \theta_ {\mathrm{F}} (^ {t} v (x ^ {*}) \otimes y)  ,
\]

and therefore

\[
\operatorname{Tr} (v \circ u) = \langle v (y), x ^ {*} \rangle = \langle y, ^ {t} v (x ^ {*}) \rangle = \operatorname{Tr} (u \circ v).
\]

COROLLARY. --- Let E be a projective A-module, u be an element of \(\operatorname{End}_{\mathrm{A}}^{f}(\mathrm{E})\) , and F be a projective A-submodule of E containing Im u. Denote by \(u_{F}\) the endomorphism of F induced by u. We have

\[
\operatorname{Tr} (u) = \operatorname{Tr} (u _ {\mathrm{F}}).\tag{4}
\]

Denote by i the canonical injection of F into E and by \(v: E \rightarrow F\) the homomorphism deduced from u. We have \(u_{F} = v \circ i\) and \(u = i \circ v\) . The corollary follows.

Let E be a projective A-module and \(u \in \operatorname{End}_{\mathrm{A}}^{\mathrm{f}}(\mathrm{E})\) . For every natural number p, the A-module \(\wedge^{p}E\) is projective (III, §7, No.~8, p.~519, Corollary 2), and the endomorphism \(\wedge^{p}u\) belongs to \(\operatorname{End}_{\mathrm{A}}^{\mathrm{f}}(\wedge^{p}\mathrm{E})\) (III, §7, No.~3, p.~511, Proposition 6) and is zero for p sufficiently large. The set \(1_{\mathrm{E}} + \operatorname{End}_{\mathrm{A}}^{\mathrm{f}}(\mathrm{E})\) is stable under composition. We define a mapping det from \(1_{\mathrm{E}} + \operatorname{End}_{\mathrm{A}}^{\mathrm{f}}(\mathrm{E})\) to A by setting

\[
\det (1 _ {\mathrm{E}} + u) = \sum_ {p \geqslant 0} \operatorname{Tr} \wedge^ {p} u
\]

for \(u\in \mathrm{End}_{\mathrm{A}}^{\mathrm{f}}(\mathrm{E})\)

If E is free and finite-dimensional, then this definition agrees with that of III, §8, No.~1, p.~522, by the corollary of III, §8, No.~5, p.~530.

PROPOSITION 3. --- Let E be a projective A-module.

\begin{enumerate}
\def\labelenumi{\alph{enumi})}
\tightlist
\item
  Let \(u \in \mathrm{End}_{\mathrm{A}}^{\mathrm{f}}(\mathrm{E})\) . Let \(\mathrm{F}\) be a projective A-submodule of \(\mathrm{E}\) containing \(\operatorname{Im} u\) , and let \(u_{\mathrm{F}}\) be the endomorphism of \(\mathrm{F}\) induced by \(u\) . We have
\end{enumerate}

\[
\det (1 _ {\mathrm{E}} + u) = \det (1 _ {\mathrm{F}} + u _ {\mathrm{F}}).\tag{5}
\]

\begin{enumerate}
\def\labelenumi{\alph{enumi})}
\setcounter{enumi}{1}
\tightlist
\item
  Let \(u\) and \(v\) be two elements of \(\mathrm{End}_{\mathrm{A}}^{\mathrm{f}}(\mathrm{E})\) . We have
\end{enumerate}

\[
\det \bigl ((1 _ {\mathrm{E}} + u) \circ (1 _ {\mathrm{E}} + v) \bigr) = \det (1 _ {\mathrm{E}} + u) \det (1 _ {\mathrm{E}} + v).\tag{6}
\]

Let us prove a). For every integer \(p \geqslant 0\) , the projective A-module \(\wedge^p\mathrm{F}\) can be identified with a submodule of \(\wedge^p\mathrm{E}\) (III, §7, No.~9, p.~520, Corollary). The image of \(\wedge^p u\) is contained in \(\wedge^p\mathrm{F}\) , and the endomorphism of \(\wedge^p\mathrm{F}\) induced by \(\wedge^{p}u\) is equal to \(\wedge^{p}u_{F}\) . We consequently have \(\operatorname{Tr}(\wedge^{p}u)=\operatorname{Tr}(\wedge^{p}u_{F})\) by the corollary of Proposition 2, and therefore a).

Let us prove b). Let G be an A-module such that the A-module \(L = E \oplus G\) is free. Denote by \(u'\) and \(v'\) the endomorphisms \(u \oplus 0_{G}\) and \(v \oplus 0_{G}\) of L. By a), we have the relations \(\det(1_{L} + u') = \det(1_{E} + u)\) , \(\det(1_{L} + v') = \det(1_{E} + v)\) , and

\[
\det (1 _ {\mathrm{L}} + u ^ {\prime} + v ^ {\prime} + u ^ {\prime} \circ v ^ {\prime}) = \det (1 _ {\mathrm{E}} + u + v + u \circ v).
\]

It therefore suffices to prove assertion b) when the A-module E is free. There then exists a finitely generated free submodule F of E that contains the image of u and that of v. Set \(w = u + v + u \circ v\) . The image of w is contained in F, and we have \(w_{F} = u_{F} + v_{F} + u_{F} \circ v_{F}\) . Therefore, by (5), we have \(\det(1_{E} + u) = \det(1_{F} + u_{F})\) , \(\det(1_{E} + v) = \det(1_{F} + v_{F})\) , and

\[
\begin{array}{c} \det \bigl ((1 _ {\mathrm{E}} + u) \circ (1 _ {\mathrm{E}} + v) \bigr) \\ = \det (1 _ {\mathrm{E}} + w) = \det (1 _ {\mathrm{F}} + w _ {\mathrm{F}}) = \det \bigl ((1 _ {\mathrm{F}} + u _ {\mathrm{F}}) \circ (1 _ {\mathrm{F}} + v _ {\mathrm{F}}) \bigr). \end{array}
\]

Since F is a finitely generated free A-module, we have

\[
\det \left(\left(1 _ {\mathrm{F}} + u _ {\mathrm{F}}\right) \circ \left(1 _ {\mathrm{F}} + v _ {\mathrm{F}}\right)\right) = \det \left(1 _ {\mathrm{F}} + u _ {\mathrm{F}}\right) \det \left(1 _ {\mathrm{F}} + v _ {\mathrm{F}}\right) = \det \left(1 _ {\mathrm{E}} + u\right) \det \left(1 _ {\mathrm{E}} + v\right).
\]

\BourbakiAppendixExercises{Exercises}

\begin{enumerate}
\def\labelenumi{\arabic{enumi})}
\tightlist
\item
  Let K be a commutative field and E a vector space over K. We denote by \(\operatorname{End}_{\mathrm{K}}^{\mathrm{pf}}(\mathrm{E})\) the subset of \(\operatorname{End}_{\mathrm{K}}(\mathrm{E})\) consisting of the endomorphisms of which a power belongs to \(\operatorname{End}_{\mathrm{K}}^{\mathrm{f}}(\mathrm{E})\) . If \(u \in \operatorname{End}_{\mathrm{K}}^{\mathrm{pf}}(\mathrm{E})\) , then we set \(I_{u} = \bigcap_{n \geqslant 0} Im u^{n}\) and \(\operatorname{Tr} u = \operatorname{Tr}(u | I_{u})\) ; this notion coincides with that defined in No.~2 for an endomorphism of finite rank. a) Let \(u \in \operatorname{End}_{\mathrm{K}}^{\mathrm{pf}}(\mathrm{E})\) ; let V be a subspace of E stable under u, and let \(u_{V}\) and \(u_{E/V}\) be the endomorphisms of V and E/V induced by u. Prove the equality \(\operatorname{Tr} u = \operatorname{Tr} u_{V} + \operatorname{Tr} u_{E/V}\) .
\end{enumerate}

\begin{enumerate}
\def\labelenumi{\alph{enumi})}
\setcounter{enumi}{1}
\item
  Let U be a subspace of \(\operatorname{End}_{\mathbf{K}}(\operatorname{E})\) . Suppose that there exists an integer n such that every product of n elements of U has finite rank. Prove that the mapping \(Tr: U \to K\) is linear.
\item
  Let F be a vector space over K, and let \(u: E \to F\) and \(v: F \to E\) be homomorphisms such that \(u \circ v \in \operatorname{End}_{\mathrm{K}}^{\mathrm{pf}}(\mathrm{F})\) . We have \(v \circ u \in \operatorname{End}_{\mathrm{K}}^{\mathrm{pf}}(\mathrm{E})\) and \(\operatorname{Tr} v \circ u = \operatorname{Tr} u \circ v\) .
\end{enumerate}

¶ 2) Keep the notation of the previous exercise. If V and W are subspaces of E, then we denote by V \(\prec\) W the relation ``W has finite codimension in V + W.''

\begin{enumerate}
\def\labelenumi{\alph{enumi})}
\item
  Let V be a subspace of E. We denote by \(\mathfrak{A}\) (resp. \(\mathfrak{a},\mathfrak{b}\) ) the set of endomorphisms \(u\) of V such that \(u(\mathrm{V})\prec \mathrm{V}\) (resp. \(u(\mathrm{E})\prec \mathrm{V}\) , \(u(\mathrm{V})\prec 0\) ). Prove that \(\mathfrak{A}\) is a K-subalgebra of \(\operatorname{End}_{\mathrm{K}}(\mathrm{E})\) , that \(\mathfrak{a}\) and \(\mathfrak{b}\) are two-sided ideals of \(\mathfrak{A}\) , with sum \(\mathfrak{A}\) , and that we have \(\mathfrak{a} \cap \mathfrak{b} \subset \operatorname{End}_{\mathrm{K}}^{\mathrm{pf}}(\mathrm{E})\) .
\item
  Let \(u, u'\) be permutable elements of \(\mathfrak{A}\) ; we can write \(u = a + b\) and \(u' = a' + b'\) with \(a, a'\) in \(\mathfrak{a}\) and \(b, b'\) in \(\mathfrak{b}\) . Prove that the element \([a, a'] = aa' - a'a\) of \(\mathfrak{A}\) belongs to \(\mathfrak{a} \cap \mathfrak{b}\) and that its trace does not depend on the chosen decompositions of \(u\) and \(u'\) ; we denote it by \(\langle u, u' \rangle\) .
\item
  Let A be a commutative subalgebra of \(\mathfrak{A}\) . Prove that there exists a unique K-linear mapping \(\operatorname{Res}:\Omega_{\mathrm{K}}(\mathrm{A})\to \mathrm{K}\) such that \(\operatorname{Res}(udv) = \langle u,v\rangle\) for all \(u,v\) in A (use the identity \([u,vw] + [v,wu] + [w,uv] = 0\) ).
\item
  Let \(u, v \in A\) ; suppose \(v(V) \subset V\) . Let \(\pi\) be a projector of E with image V. Prove the equality \(\operatorname{Res}(u dv) = \operatorname{Tr}[\pi u, v]\) . Deduce that if \(u(V) \subset V\) , then we have \(\operatorname{Res} u dv = 0\) (observe that the endomorphism \([\pi u, v]\) has square zero).
\item
  Let \(u \in \mathrm{A}\) , \(n \in \mathbf{N}\) . Prove that we have \(\operatorname{Res} u^n du = 0\) . If \(u\) is invertible in \(\mathrm{A}\) , then we have \(\operatorname{Res} u^{-n} du = 0\) for \(n \geqslant 2\) ; if moreover \(u(\mathrm{V}) \subset \mathrm{V}\) , then we have \(\operatorname{Res} u^{-1} du = \dim_{\mathrm{K}} \mathrm{V} / u(\mathrm{V})\) .
\end{enumerate}

\(f)\) Take \(\mathrm{E} = \mathrm{K}((\mathrm{X}))\) (IV, §4, No.~9, p.~38), \(\mathrm{V} = \mathrm{K}[[\mathrm{X}]]\) , and let \(\mathrm{A} = \mathrm{E}\) act on itself by homotheties. Prove that the assumption of \(c)\) is satisfied. Let \(f(\mathrm{X}) = \sum_{n} a_{n} \mathrm{X}^{n} d\mathrm{X}\) be an element of \(\mathrm{K}((\mathrm{X}))\) . Prove that \(\operatorname{Res} f(\mathrm{X}) d\mathrm{X}\) is equal to \(a_{-1}\) . Deduce that if \(g\) is an element of \(\mathrm{K}[[\mathrm{X}]]\) such that \(g(0) = 0\) and \(g'(0) \neq 0\) , then the coefficients of \(\mathrm{X}^{-1}\) in the formal series expansions of \(f(\mathrm{X})\) and \(f(g(\mathrm{X})) g'( \mathrm{X})\) are the same.

\backmatter
\BourbakiBackMatter{Historical Note}

We have seen (historical notes of Chapters I and II--III) that the first noncommutative algebras appear in 1843--44, in works of Hamilton {[}25{]} and Grassmann {[}23{]} and {[}24{]}. Hamilton, when introducing the quaternions, already has a clear idea of arbitrary algebras of finite rank over the field of real numbers ({[}25{]}, Preface, p.~26--31) \(^{(1)}\) . In developing his theory, he later has the idea of considering what he calls ``biquaternions,'' that is, the algebra over the field of complex numbers with the same multiplication table as the quaternion algebra; on this occasion, he observes that this extension causes the appearance of zero divisors ({[}25{]}, p.~650). Grassmann's point of view is slightly different, and for a long time, his ``exterior algebra'' would remain separate from the general theory of algebras; \(^{(2)}\) but in his language, which still lacks precision, one cannot fail to recognize the first idea of an algebra (of finite dimension or not, over the field of real numbers) defined by a generating set and relations {[}24{]}.

New examples of algebras are introduced, more or less explicitly, between 1850 and 1860. Cayley, when developing the theory of matrices {[}8{]}, does not yet see the square matrices as forming an algebra (a point of view that will first be clearly expressed by the Peirces around 1870 {[}40{]}). However, he does already note, on this occasion, the existence of a system of matrices of order 2 satisfying the multiplication table of the quaternions; this observation can be seen as the first example of a linear representation of an algebra. \(^{(3)}\) Moreover, in the memoir where he defines the abstract notion of finite group, he also gives, in passing, the definition of the algebra of such a group, without further applying the definition ({[}7{]}, p.~129).

No other progress is made before 1870, but at this time, research into the general structure of finite-dimensional algebras (over the field of real or complex numbers) begins. It is B. Peirce who takes the first steps in this direction; he introduces the notions of nilpotent element and of idempotent element, proves that an algebra (with or without unit) with at least one nonnilpotent element has a nonzero idempotent, writes down the famous decomposition

\[
x = e x e + (x e - e x e) + (e x - e x e) + (x - x e - e x + e x e)
\]

(e an idempotent, x an arbitrary element), and has the (still slightly imprecise) idea of a decomposition of an idempotent as a sum of pairwise orthogonal ``primitive'' idempotents {[}40{]}. Moreover, according to Clifford ({[}11{]}, p.~274), \(^{(4)}\) it is to B. Pierce that should be attributed the notion of tensor product of two algebras, which Clifford himself applies implicitly to a generalization of Hamilton's ``biquaternions'' {[}9{]} and explicitly to the study of the algebras named after him, several years later ({[}10{]} and {[}11{]}). These new notions are used by B. Peirce to classify low-dimensional algebras (over the field of complex numbers). This problem was also tackled, around 1880, by other mathematicians of the Anglo-American school, in particular Cayley and Sylvester. The great variety of possible structures quickly becomes apparent, and it is undoubtedly this fact that will, in the following period, direct research toward obtaining classes of algebras with more particular properties.

In continental Europe, where the ideas evolve quite differently, such research already appears before 1880. In 1878, Frobenius proves that the quaternions form the only example of a noncommutative (finite-dimensional) field over the field of real numbers ({[}18{]}, p.~59--63), a result published independently two years later by C. S. Pierce {[}39{]}. As early as 1861, Weierstrass, in clarifying a remark by Gauss, had, in his lectures, characterized commutative algebras over R or C without nilpotent elements \(^{(5)}\) as direct sums of fields isomorphic to R or C. Dedekind, in his research, had reached the same conclusions around 1870, in connection with his ``hypercomplex'' conception of the theory of commutative fields. Their proofs are published in 1884--85 ({[}54{]} and {[}12{]}). It is also in 1884 that H. Poincaré, in a quite cryptic note {[}41{]}, draws attention to the possibility of viewing the equations \(z_{i} = \varphi_{i}(x_{1}, \ldots, x_{n}, y_{1}, \ldots, y_{n})\) that express the multiplicative law \((\sum_{i} x_{i} e_{i})(\sum_{i} y_{i} e_{i}) = \sum_{i} z_{i} e_{i}\) in an algebra as defining (locally, of course) a Lie group. This remark seems to have made a great impression on Lie and his followers (E. Study, G. Scheffers, F. Schur, and a bit later T. Molien and É. Cartan) who, at that time, were developing the theory of ``continuous'' groups and especially classification problems (see, in particular, {[}43{]}, p.~387). During the period 1885--1905, it leads the mathematicians of this school to apply methods similar to those they use when studying Lie groups and Lie algebras to the study of the structure of algebras. These methods are based, above all, on the consideration of the characteristic polynomial of an element of the algebra with respect to its regular representation (a polynomial already encountered in the works of Weierstrass and Dedekind mentioned above) and on the decomposition of this polynomial into irreducible factors. That decomposition, as Frobenius would discover a little later, is reflected in the decomposition of the regular representation into irreducible components.

In the course of the research carried out by Lie's school on algebras, the ``intrinsic'' notions of the theory gradually emerge. The notion of radical appears in a specific case (that when the quotient by the radical is a direct product of fields) in G. Scheffers' work in 1891 {[}43{]}, and more clearly in that of T. Molien {[}31{]} and É. Cartan {[}4{]}, who study the general case (the word radical comes from Frobenius {[}22{]}). Study and Scheffers {[}43{]} highlight the concept of an algebra that is a direct product of several others (already glimpsed by Peirce ({[}38{]}, p.~221). Finally, Molien {[}31{]} introduces quotient algebras of an algebra, a notion essentially equivalent to that of two-sided ideal (first defined by É. Cartan {[}4{]}) or of homomorphism (a name also due to Frobenius). The analogy with groups is obvious, and later, in 1904, Epsteen and Wedderburn will consider composition series of two-sided ideals and extend the Jordan--Hölder theorem to them. The most important results of this period are T. Molien's {[}31{]}: guided by the notion of simple group, he defines simple algebras (over C) and proves that these are matrix algebras. He then proves that the structure of an arbitrary algebra of finite rank over C essentially reduces to the case, already studied by Scheffers, when the quotient by the radical is a direct sum of fields. Shortly after that, these results are rediscovered and established more rigorously and more clearly by É. Cartan {[}4{]}. On this occasion, he introduces the notion of semisimple algebra and brings to light numerical invariants, the ``É. Cartan integers'' (VIII, p.~180, Exercise 8), associated with an arbitrary algebra over the field C. He thus brings the theory of these algebras to a point beyond which little progress has been made since. \(^{(6)}\) Finally, he extends Molien's results and his own on algebras over R.

Around 1900, ideas begin to develop in a direction that leads to abandoning every restriction on the field of scalars in everything related to linear algebra. In particular, we must point out the vigorous impulse given to studying finite fields by the American school surrounding E. H. Moore and L. E. Dickson. The result of this research is Wedderburn's theorem {[}50{]} that proves that every finite field is commutative. In 1907, Wedderburn takes Cartan's results and extends them to an arbitrary base field {[}51{]}. In doing so, he completely abandons his predecessors' methods (that cannot be applied when the field is no longer algebraically closed or maximal ordered). Instead, he returns, while refining it, to B. Peirce's technique of idempotents, which allows him to put in definite form the structure theorem for semisimple algebras, reducing their study to that of noncommutative fields. Moreover, the problems of the extension of scalars arise naturally in his perspective, and Wedderburn proves that every semisimple algebra remains semisimple after a separable extension of the base field \(^{(7)}\) and becomes a direct product of central simple matrix algebras when this extension is sufficiently large ({[}51{]}, p.~102). \(^{(8)}\) A bit later, Dickson, for n = 3 {[}17{]}, and Wedderburn himself, for n arbitrary {[}52{]}, give the first examples of noncommutative fields of rank \(n^{2}\) over their center, \(^{(9)}\) thus introducing in a specific case the theory of ``cross products'' and of ``systems of factors'' that would later be developed by R. Brauer {[}2{]} and E. Noether ({[}33{]}, {[}34{]}, and {[}35{]}). Finally, in 1921, Wedderburn proves a specific case of the commutation theorem {[}53{]}.

In the meantime, from 1896 to 1910, a theory close to that of algebras, the theory of linear representations of groups (limited at first to representations of finite groups), was developed by Frobenius, Burnside, and I. Schur. It finds its origin in remarks made by Dedekind. In the course of his research on normal bases of Galois extensions, Dedekind had (even before the publication of his work on algebras), around 1880, come across the ``Gruppendeterminant'' \(\det(x_{st^{-1}})\) , where \((x_{s})_{s\in G}\) is a sequence of variables whose index set is a finite group G (in other words, the norm of the generic element of the algebra of the group G with respect to its regular representation). He had observed that when G is abelian, this polynomial decomposes into linear factors (which generalizes an identity proved long before for determinants of ``circulant'' matrices, which correspond to cyclic groups G). In the course of his fascinating correspondence with Frobenius {[}13{]}, Dedekind, in 1896, draws his attention to this property, its link to the theory of characters of abelian groups, and analogous results on specific noncommutative groups that he obtained in 1886. A few months later, Frobenius completely solves the problem of the decomposition of the ``Gruppendeterminant'' into irreducible factors {[}19{]}, thanks to his brilliant generalization of the notion of character {[}20{]}, which we will not discuss here. However, it should be noted that in the later development of this theory, \(^{(10)}\) Frobenius always remains aware of its connection to the theory of algebras (on which Dedekind had not ceased to insist in his letters). After having introduced the notions of irreducible representation and of completely reducible representation for groups {[}21{]} and proved that the regular representation contains all irreducible representations, it is through analogous methods that he proposes, in 1903, to take up the theory of Molien--Cartan {[}22{]}. In the work of Burnside {[}3{]} and I. Schur {[}45{]}, the ``hypercomplex'' aspect of the theory does not explicitly intervene; but it is in their work that the fundamental properties of irreducible representations, Schur's lemma and Burnside's theorem, appear. Finally, it should be noted that it is in this theory that two particular cases of the commutation theorem appear for the first time. The first is in I. Schur's thesis {[}44{]}, which links (precisely by commutation in the endomorphism ring of a tensor space) the representations of the linear group and those of the symmetric group. The second is in Schur's 1905 work {[}46{]}, where he shows that matrices that commute with all matrices of an irreducible representation over the field C are scalar multiples of \(I\) (a result that also follows from Burnside's theorem).

It remained to clearly identify the common base of these theories: this was the work of the German school surrounding E. Noether and E. Artin, in the period 1925--1933, which sees the creation of modern algebra. Already in 1903, in a memoir on the algebraic integration of linear differential equations {[}42{]}, H. Poincaré had defined, in an algebra, left and right ideals and the notion of minimal ideal. He had also observed that in a semisimple algebra, every left ideal is the direct sum of its intersections with the simple components and that in the algebra of matrices of order n, the minimal ideals have dimension n; but his work went unnoticed by algebraists. \(^{(11)}\) In 1907, Wedderburn defines left and right ideals of an algebra again and proves some of their properties (in particular, that the radical is the largest nilpotent left ideal ({[}51{]}, p.~113--114)).

But it is not until 1927 that these notions are used in an essential way in the theory of algebras. \(^{(12)}\) Generalizing proof methods that had appeared here and there, \(^{(13)}\) W. Krull in 1925 {[}28{]} and E. Noether in 1926 {[}33{]} introduce and systematically use the maximal and minimal conditions. The first uses them to extend Remak's theorem on the decomposition of a finite group as a direct product of indecomposable groups to abelian groups with operators (that he defines on this occasion) (cf.~VIII, p.~37, Theorem 2), while the second uses these conditions in the characterization of Dedekind rings. In 1927, E. Artin {[}1{]}, applying the same idea to noncommutative rings, shows how, through a systematic study of minimal ideals, one can extend Wedderburn's theorem to all rings whose left ideals satisfy both the maximal and the minimal conditions. \(^{(14)}\)

Elsewhere, in 1926, Krull {[}29{]} makes the link between the notion of abelian group with operators and that of linear representation of a group. This point of view is generalized to algebras and developed in detail by E. Noether in a fundamental 1929 work {[}34{]}. Because of the importance of the introduced ideas and the clarity of the exposition, that work deserves to be listed next to Steinitz' memoir on commutative fields as one of the pillars of modern algebra. \(^{(15)}\)

Finally, in a series of works beginning in 1927 ({[}36{]}, {[}2{]}, and {[}35{]}), E. Noether and R. Brauer, joined in 1929 by A. Albert and H. Hasse, take up the study of skew fields where Wedderburn and Dickson left off. Although the most essential part of their results consists in a thorough study of the Brauer group, in particular over algebraic number fields, it is in these works that are clarified the commutation theorems as well as the notion of neutralizing field of a simple algebra and its relation to maximal commutative subfields. Finally, in 1927, Skolem characterizes the automorphisms of simple rings {[}49{]} in a theorem recovered several years later by E. Noether {[}35{]} and R. Brauer {[}2{]}.

Thus, in 1934, the elementary theory of simple and semisimple rings has more or less reached its final form (for an overview of the current state of the theory, see {[}16{]}).
